% ============================================================
% Sección 8 — Malla curricular por grado y periodo
% Generado automáticamente por gen_tablas.py
% ============================================================

\subsubsection{Periodo 2 · Hardware, software y la máquina por dentro}
\noindent\begin{tabular}{ll}
\textbf{Grado:} Sexto & \textbf{Periodo:} Segundo \\
\textbf{Proyecto integrador:} Maestros del oficio digital: abrir la máquina & \textbf{Anclaje:} Relojero, costurera y panadero de Cartago \\
\end{tabular}
\medskip

\begin{longtable}{|L{3.8cm}|L{3.3cm}|L{3.3cm}|L{2.5cm}|L{2.5cm}|L{2.5cm}|}
\caption{Estructura de malla curricular — G6° Periodo 2 · Hardware, software y la máquina por dentro}
\label{tab:g6p2}\\
\hline
\rowcolor{azulInst}
\multirow{2}{*}{\color{white}\scriptsize\textbf{\makecell{ESTÁNDARES, LINEAMIENTOS\\U ORIENTACIONES}}} &
\multirow{2}{*}{\color{white}\scriptsize\textbf{\makecell{COMPETENCIAS}}} &
\multirow{2}{*}{\color{white}\scriptsize\textbf{\makecell{APRENDIZAJES}}} &
\multicolumn{3}{c|}{%
  \color{white}\makecell{\textbf{\footnotesize DESCRIPTORES DE DESEMPEÑO}\\[1pt]\textit{\scriptsize Metas del aprendizaje por periodo}}%
}\\
\cline{4-6}
\rowcolor{azulClaro}
\cellcolor{azulInst} & \cellcolor{azulInst} & \cellcolor{azulInst} &
\textbf{\scriptsize COGNITIVOS} &
\textbf{\scriptsize PROCEDIMENTALES} &
\textbf{\scriptsize ACTITUDINALES} \\
\hline
\endfirsthead
\hline
\rowcolor{azulInst}
\multirow{2}{*}{\color{white}\scriptsize\textbf{\makecell{ESTÁNDARES, LINEAMIENTOS\\U ORIENTACIONES}}} &
\multirow{2}{*}{\color{white}\scriptsize\textbf{\makecell{COMPETENCIAS}}} &
\multirow{2}{*}{\color{white}\scriptsize\textbf{\makecell{APRENDIZAJES}}} &
\multicolumn{3}{c|}{%
  \color{white}\makecell{\textbf{\footnotesize DESCRIPTORES DE DESEMPEÑO}\\[1pt]\textit{\scriptsize Metas del aprendizaje por periodo}}%
}\\
\cline{4-6}
\rowcolor{azulClaro}
\cellcolor{azulInst} & \cellcolor{azulInst} & \cellcolor{azulInst} &
\textbf{\scriptsize COGNITIVOS} &
\textbf{\scriptsize PROCEDIMENTALES} &
\textbf{\scriptsize ACTITUDINALES} \\
\hline
\endhead

% -- FILA 1: Naturaleza y evolución --
\textbf{Naturaleza y evolución de la tecnología} &
Comprender las 4 funciones básicas de toda máquina inteligente y reconocer la evolución del hardware. &
Identifica las 4 funciones (E/P/S/A) y las piezas del gabinete (S2–S3). &
Reconoce CPU, RAM, disco, fuente y tarjeta madre, comprendiendo el papel de cada pieza en el funcionamiento del sistema. &
Elabora un mapa anatómico del computador con 8 piezas identificadas y descripción propia. &
Valora la habilidad del relojero/costurera/panadero como antecedente del técnico actual. \\
\hline

% -- FILA 2: Apropiación y uso --
\textbf{Apropiación y uso de la tecnología} &
Manejar periféricos, sistema operativo y archivos digitales con criterio de organización profesional. &
Distingue periféricos de entrada, salida y mixtos (S4), comprende qué es el SO (S6) y organiza archivos en carpetas (S7). &
Diferencia hardware de software y conoce los principales sistemas operativos. &
Organiza sus archivos en estructura profesional (año/materia/periodo) con nomenclatura correcta. &
Aplica buenas prácticas en la sala de sistemas y respeta los equipos compartidos. \\
\hline

% -- FILA 3: Solución de problemas --
\textbf{Solución de problemas con tecnología} &
Diagnosticar problemas básicos de un equipo informático y aplicar mantenimiento preventivo. &
Aplica las 5 reglas de cuidado de la sala (S1, S8, S9) y reconoce señales de fallo (sobrecalentamiento, lentitud). &
Identifica las 3 acciones de mantenimiento básico: limpiar, actualizar, respaldar. &
Realiza limpieza física suave del equipo, actualiza el SO y crea respaldos de archivos. &
Cuida su cuerpo aplicando la regla 20-20-20 y la postura ergonómica. \\
\hline

% -- FILA 4: Tecnología y sociedad --
\textbf{Tecnología y sociedad} &
Elegir un equipo informático con criterio razonado, evitando el extremismo del precio. &
Construye una ficha técnica del computador ideal para estudio con presupuesto razonado (S10). &
Reconoce las especificaciones mínimas razonables para uso escolar (CPU i3/Ryzen3, RAM 8 GB, SSD 256 GB). &
Elabora una ficha técnica con marca, modelo, especificaciones y precio realista. &
Aplica criterio frente a la publicidad; reconoce que lo barato sale caro. \\
\hline

% -- FILA 5: PC (verde) --
\rowcolor{verdePCclaro}
\makecell[l]{\textbf{\color{verdePC}Pensamiento}\\\textbf{\color{verdePC}Computacional}\\
\footnotesize Wing (2006)\\\footnotesize MEN Or. 2022\\\footnotesize\textit{Descomposición · Abstracción}}\
 &
Aplicar la descomposición y la abstracción para comprender el sistema informático como capas funcionales interrelacionadas. &
Descompone el computador en sus 4 funciones (E/P/S/A) y abstrae el SO como capa de intermediación (S2–S6). &
Describe el computador como sistema jerárquico de capas (hardware / SO / aplicaciones / usuario). &
Elabora un diagrama de capas del sistema informático con funciones interrelacionadas. &
Valora la arquitectura del computador como diseño humano inteligente; entiende el sistema para usarlo con criterio. \\
\hline

% -- FILA 6: SEL (violeta) --
\rowcolor{moradoSELclaro}
\makecell[l]{\textbf{\color{moradoSEL}Habilidades}\\\textbf{\color{moradoSEL}Socioemocionales}\\
\footnotesize Ley 2383/2024\\\footnotesize CASEL 2020\\\footnotesize\textit{Autogestión · Conciencia social}}\
 &
Desarrollar autogestión de la frustración frente a equipos que fallan y conciencia del cuidado colectivo de los bienes tecnológicos compartidos. &
Identifica sus propias reacciones emocionales cuando un equipo falla y reconoce que los bienes tecnológicos de la sala son responsabilidad colectiva. &
Distingue entre frustración individual y responsabilidad colectiva en el uso de equipos compartidos. &
Aplica la regla de pausa-respira-reporta cuando un equipo presenta fallas, sin ocultarlas. &
Cuida los equipos de la sala como bienes comunes y reporta fallas con honestidad para el bien del grupo. \\
\hline

\end{longtable}
\medskip
\noindent\textbf{Nota:} Las filas 5 (\textcolor{verdePC}{verde}) y 6 (\textcolor{moradoSEL}{violeta}) son componentes propios de este plan, adicionales a los cuatro componentes de la Guía No.~30 del MEN.

\subsubsection{Periodo 3 · Escribir y buscar con criterio}
\noindent\begin{tabular}{ll}
\textbf{Grado:} Sexto & \textbf{Periodo:} Tercero \\
\textbf{Proyecto integrador:} Reporteros con criterio del barrio & \textbf{Anclaje:} Doña Mercedes, maestra rural (La Plata, Cartago) \\
\end{tabular}
\medskip

\begin{longtable}{|L{3.8cm}|L{3.3cm}|L{3.3cm}|L{2.5cm}|L{2.5cm}|L{2.5cm}|}
\caption{Estructura de malla curricular — G6° Periodo 3 · Escribir y buscar con criterio}
\label{tab:g6p3}\\
\hline
\rowcolor{azulInst}
\multirow{2}{*}{\color{white}\scriptsize\textbf{\makecell{ESTÁNDARES, LINEAMIENTOS\\U ORIENTACIONES}}} &
\multirow{2}{*}{\color{white}\scriptsize\textbf{\makecell{COMPETENCIAS}}} &
\multirow{2}{*}{\color{white}\scriptsize\textbf{\makecell{APRENDIZAJES}}} &
\multicolumn{3}{c|}{%
  \color{white}\makecell{\textbf{\footnotesize DESCRIPTORES DE DESEMPEÑO}\\[1pt]\textit{\scriptsize Metas del aprendizaje por periodo}}%
}\\
\cline{4-6}
\rowcolor{azulClaro}
\cellcolor{azulInst} & \cellcolor{azulInst} & \cellcolor{azulInst} &
\textbf{\scriptsize COGNITIVOS} &
\textbf{\scriptsize PROCEDIMENTALES} &
\textbf{\scriptsize ACTITUDINALES} \\
\hline
\endfirsthead
\hline
\rowcolor{azulInst}
\multirow{2}{*}{\color{white}\scriptsize\textbf{\makecell{ESTÁNDARES, LINEAMIENTOS\\U ORIENTACIONES}}} &
\multirow{2}{*}{\color{white}\scriptsize\textbf{\makecell{COMPETENCIAS}}} &
\multirow{2}{*}{\color{white}\scriptsize\textbf{\makecell{APRENDIZAJES}}} &
\multicolumn{3}{c|}{%
  \color{white}\makecell{\textbf{\footnotesize DESCRIPTORES DE DESEMPEÑO}\\[1pt]\textit{\scriptsize Metas del aprendizaje por periodo}}%
}\\
\cline{4-6}
\rowcolor{azulClaro}
\cellcolor{azulInst} & \cellcolor{azulInst} & \cellcolor{azulInst} &
\textbf{\scriptsize COGNITIVOS} &
\textbf{\scriptsize PROCEDIMENTALES} &
\textbf{\scriptsize ACTITUDINALES} \\
\hline
\endhead

% -- FILA 1: Naturaleza y evolución --
\textbf{Naturaleza y evolución de la tecnología} &
Comprender qué es internet, cómo viajan los datos y la diferencia entre internet y la World Wide Web. &
Reconoce la arquitectura cliente-servidor y los componentes de internet (S7). &
Describe internet como red de redes conectada por cables submarinos, satélites y antenas. &
Explica con ejemplos cómo viaja una petición desde su navegador hasta un servidor y de regreso. &
Aprecia internet como bien común global que requiere cuidado y mantenimiento colectivo. \\
\hline

% -- FILA 2: Apropiación y uso --
\textbf{Apropiación y uso de la tecnología} &
Producir documentos digitales con formato profesional usando un procesador de texto. &
Aplica formato básico, listas, tablas, imágenes y encabezados profesionales (S2–S6). &
Identifica los elementos de un documento profesional: fuente, tamaño, márgenes, alineación, listas, tablas, encabezado, pie y números de página. &
Produce un documento informativo de 4 páginas con todos los elementos de formato profesional. &
Cuida la presentación y valora la legibilidad como respeto al lector. \\
\hline

% -- FILA 3: Solución de problemas --
\textbf{Solución de problemas con tecnología} &
Buscar información con operadores avanzados y evaluar fuentes con criterio CRAAP para detectar fake news. &
Aplica operadores (site:, comillas, guion menos) y evalúa fuentes con las 5 letras CRAAP (S8–S9). &
Reconoce los 5 criterios CRAAP y las 5 señales típicas de fake news. &
Investiga un tema usando al menos 3 fuentes evaluadas con CRAAP; elabora guía visual anti-fake news. &
Antes de creer o reenviar, se detiene, respira y verifica; no comparte rumores. \\
\hline

% -- FILA 4: Tecnología y sociedad --
\textbf{Tecnología y sociedad} &
Producir y publicar información con responsabilidad ciudadana, citando fuentes y respetando autoría. &
Cita fuentes correctamente, evita el plagio y reconoce la responsabilidad de quien publica (S10). &
Identifica la importancia de citar fuentes, la honestidad académica y el rol del productor responsable en la infoesfera. &
Entrega un informe ciudadano del barrio con mínimo 2 fuentes citadas y declaración de uso de IA. &
Asume su rol como ciudadano digital productor y aporta claridad a su comunidad. \\
\hline

% -- FILA 5: PC (verde) --
\rowcolor{verdePCclaro}
\makecell[l]{\textbf{\color{verdePC}Pensamiento}\\\textbf{\color{verdePC}Computacional}\\
\footnotesize Wing (2006)\\\footnotesize MEN Or. 2022\\\footnotesize\textit{Diseño de algoritmos · Evaluación}}\
 &
Aplicar algoritmos de búsqueda y evaluación de fuentes como procesos lógicos secuenciales para tomar decisiones informadas. &
Aplica operadores de búsqueda como instrucciones secuenciales; usa CRAAP como algoritmo de verificación de 5 pasos. &
Reconoce los operadores de búsqueda como instrucciones precisas para un motor algorítmico; describe CRAAP como árbol de decisión. &
Diseña su propio flujo de búsqueda y evaluación como diagrama de pasos con condiciones visibles. &
Rechaza el pensamiento impulsivo y adopta la lógica de quien verifica antes de actuar. \\
\hline

% -- FILA 6: SEL (violeta) --
\rowcolor{moradoSELclaro}
\makecell[l]{\textbf{\color{moradoSEL}Habilidades}\\\textbf{\color{moradoSEL}Socioemocionales}\\
\footnotesize Ley 2383/2024\\\footnotesize CASEL 2020\\\footnotesize\textit{Autogestión · Toma de decisiones responsable}}\
 &
Regular el impulso de creer y compartir sin verificar, y tomar decisiones responsables como productor de información en la infoesfera. &
Identifica la emoción de urgencia que genera el contenido viral y reconoce cuándo esa emoción puede llevar a errores de juicio. &
Describe los efectos sociales de compartir información falsa y el papel del ciudadano digital en frenarlo. &
Aplica la secuencia pausa-verifica-decide antes de compartir cualquier contenido en redes. &
Asume responsabilidad por la información que produce y difunde, reconociendo su impacto en la comunidad. \\
\hline

\end{longtable}
\medskip
\noindent\textbf{Nota:} Las filas 5 (\textcolor{verdePC}{verde}) y 6 (\textcolor{moradoSEL}{violeta}) son componentes propios de este plan, adicionales a los cuatro componentes de la Guía No.~30 del MEN.

\subsection{Grado Séptimo}

\subsubsection{Periodo 1 · Trabajo colaborativo en la nube (Microsoft 365)}
\noindent\begin{tabular}{ll}
\textbf{Grado:} Séptimo & \textbf{Periodo:} Primero \\
\textbf{Proyecto integrador:} Minga digital: una obra colaborativa en la nube & \textbf{Anclaje:} La minga del Valle del Cauca \\
\end{tabular}
\medskip

\begin{longtable}{|L{3.8cm}|L{3.3cm}|L{3.3cm}|L{2.5cm}|L{2.5cm}|L{2.5cm}|}
\caption{Estructura de malla curricular — G7° Periodo 1 · Trabajo colaborativo en la nube (Microsoft 365)}
\label{tab:g7p1}\\
\hline
\rowcolor{azulInst}
\multirow{2}{*}{\color{white}\scriptsize\textbf{\makecell{ESTÁNDARES, LINEAMIENTOS\\U ORIENTACIONES}}} &
\multirow{2}{*}{\color{white}\scriptsize\textbf{\makecell{COMPETENCIAS}}} &
\multirow{2}{*}{\color{white}\scriptsize\textbf{\makecell{APRENDIZAJES}}} &
\multicolumn{3}{c|}{%
  \color{white}\makecell{\textbf{\footnotesize DESCRIPTORES DE DESEMPEÑO}\\[1pt]\textit{\scriptsize Metas del aprendizaje por periodo}}%
}\\
\cline{4-6}
\rowcolor{azulClaro}
\cellcolor{azulInst} & \cellcolor{azulInst} & \cellcolor{azulInst} &
\textbf{\scriptsize COGNITIVOS} &
\textbf{\scriptsize PROCEDIMENTALES} &
\textbf{\scriptsize ACTITUDINALES} \\
\hline
\endfirsthead
\hline
\rowcolor{azulInst}
\multirow{2}{*}{\color{white}\scriptsize\textbf{\makecell{ESTÁNDARES, LINEAMIENTOS\\U ORIENTACIONES}}} &
\multirow{2}{*}{\color{white}\scriptsize\textbf{\makecell{COMPETENCIAS}}} &
\multirow{2}{*}{\color{white}\scriptsize\textbf{\makecell{APRENDIZAJES}}} &
\multicolumn{3}{c|}{%
  \color{white}\makecell{\textbf{\footnotesize DESCRIPTORES DE DESEMPEÑO}\\[1pt]\textit{\scriptsize Metas del aprendizaje por periodo}}%
}\\
\cline{4-6}
\rowcolor{azulClaro}
\cellcolor{azulInst} & \cellcolor{azulInst} & \cellcolor{azulInst} &
\textbf{\scriptsize COGNITIVOS} &
\textbf{\scriptsize PROCEDIMENTALES} &
\textbf{\scriptsize ACTITUDINALES} \\
\hline
\endhead

% -- FILA 1: Naturaleza y evolución --
\textbf{Naturaleza y evolución de la tecnología} &
Comprender el paso de Office tradicional a Microsoft 365 y el concepto de computación en la nube. &
Identifica qué es la nube, diferencia Office tradicional de M365 y conoce OneDrive (S2–S3). &
Describe la arquitectura cliente-nube y el rol del navegador como portal universal de acceso. &
Accede a office.com con su cuenta institucional y explora OneDrive, Word, Excel y PowerPoint en línea. &
Reconoce la minga del Valle como ancestro cultural de la coautoría digital. \\
\hline

% -- FILA 2: Apropiación y uso --
\textbf{Apropiación y uso de la tecnología} &
Trabajar en equipo en archivos compartidos con coautoría asincrónica y comentarios constructivos. &
Comparte archivos con permisos, edita en coautoría asincrónica y usa comentarios (S5–S7). &
Distingue los 3 niveles de permisos (lectura, comentarios, edición) y comprende cuándo usar cada uno. &
Coedita un Word de 4 páginas con 2 compañeros, usa comentarios para sugerencias y resuelve revisiones. &
Respeta la autoría del compañero, hace comentarios constructivos y participa con reciprocidad como en la minga. \\
\hline

% -- FILA 3: Solución de problemas --
\textbf{Solución de problemas con tecnología} &
Recuperar versiones anteriores de un archivo y gestionar conflictos de edición con criterio. &
Aplica el historial de versiones cuando algo se daña en el archivo compartido (S8). &
Identifica que Microsoft 365 guarda automáticamente versiones y permite recuperar contenido perdido. &
Practica restaurar versiones anteriores cuando algo se daña, sin pánico. &
Asume la solución técnica y humana frente a errores del equipo. \\
\hline

% -- FILA 4: Tecnología y sociedad --
\textbf{Tecnología y sociedad} &
Gestionar reuniones virtuales y comunicación profesional con Outlook y Teams. &
Maneja Outlook y Teams en el contexto escolar (S9). &
Reconoce los protocolos de invitación, aceptación y cortesía en reuniones virtuales. &
Acepta invitaciones por Outlook, se une a reuniones por Teams y participa con criterio profesional. &
Aprecia la nube como bien común del equipo: cuida lo que pone allí porque es ambiente compartido. \\
\hline

% -- FILA 5: PC (verde) --
\rowcolor{verdePCclaro}
\makecell[l]{\textbf{\color{verdePC}Pensamiento}\\\textbf{\color{verdePC}Computacional}\\
\footnotesize Wing (2006)\\\footnotesize MEN Or. 2022\\\footnotesize\textit{Abstracción · Reconocimiento de patrones}}\
 &
Aplicar la abstracción para modelar procesos colaborativos en la nube y comprender la sincronización como algoritmo distribuido. &
Abstrae el flujo de coedición como proceso con reglas de concurrencia (S5–S7); modela conflictos de versiones como problema algorítmico. &
Describe la nube como sistema distribuido con reglas de sincronización, permisos y estados. &
Elabora un diagrama de flujo del proceso de coedición: estados (editando/guardado/en conflicto) y transiciones. &
Comprende que los sistemas digitales tienen lógica interna; actúa con criterio sobre los estados y permisos. \\
\hline

% -- FILA 6: SEL (violeta) --
\rowcolor{moradoSELclaro}
\makecell[l]{\textbf{\color{moradoSEL}Habilidades}\\\textbf{\color{moradoSEL}Socioemocionales}\\
\footnotesize Ley 2383/2024\\\footnotesize CASEL 2020\\\footnotesize\textit{Habilidades relacionales · Conciencia social}}\
 &
Desarrollar habilidades de colaboración auténtica en equipos digitales, respetando la autoría y aportando con reciprocidad como en la minga del Valle. &
Identifica las normas de colaboración digital justa y reconoce cómo sus acciones en el archivo compartido afectan al equipo. &
Describe la reciprocidad de la minga como modelo de trabajo colectivo aplicable a la coautoría digital. &
Participa activamente en la coedición del documento, hace aportes concretos y usa comentarios para dar retroalimentación respetuosa. &
Cuida la autoría de sus compañeros como cuida la propia; actúa con la reciprocidad de quien trabaja en minga. \\
\hline

\end{longtable}
\medskip
\noindent\textbf{Nota:} Las filas 5 (\textcolor{verdePC}{verde}) y 6 (\textcolor{moradoSEL}{violeta}) son componentes propios de este plan, adicionales a los cuatro componentes de la Guía No.~30 del MEN.

\subsubsection{Periodo 2 · Algoritmia y pensamiento computacional (Scratch)}
\noindent\begin{tabular}{ll}
\textbf{Grado:} Séptimo & \textbf{Periodo:} Segundo \\
\textbf{Proyecto integrador:} Tejedoras digitales: mi primer programa con criterio & \textbf{Anclaje:} Doña Sofía, la tejedora del Valle \\
\end{tabular}
\medskip

\begin{longtable}{|L{3.8cm}|L{3.3cm}|L{3.3cm}|L{2.5cm}|L{2.5cm}|L{2.5cm}|}
\caption{Estructura de malla curricular — G7° Periodo 2 · Algoritmia y pensamiento computacional (Scratch)}
\label{tab:g7p2}\\
\hline
\rowcolor{azulInst}
\multirow{2}{*}{\color{white}\scriptsize\textbf{\makecell{ESTÁNDARES, LINEAMIENTOS\\U ORIENTACIONES}}} &
\multirow{2}{*}{\color{white}\scriptsize\textbf{\makecell{COMPETENCIAS}}} &
\multirow{2}{*}{\color{white}\scriptsize\textbf{\makecell{APRENDIZAJES}}} &
\multicolumn{3}{c|}{%
  \color{white}\makecell{\textbf{\footnotesize DESCRIPTORES DE DESEMPEÑO}\\[1pt]\textit{\scriptsize Metas del aprendizaje por periodo}}%
}\\
\cline{4-6}
\rowcolor{azulClaro}
\cellcolor{azulInst} & \cellcolor{azulInst} & \cellcolor{azulInst} &
\textbf{\scriptsize COGNITIVOS} &
\textbf{\scriptsize PROCEDIMENTALES} &
\textbf{\scriptsize ACTITUDINALES} \\
\hline
\endfirsthead
\hline
\rowcolor{azulInst}
\multirow{2}{*}{\color{white}\scriptsize\textbf{\makecell{ESTÁNDARES, LINEAMIENTOS\\U ORIENTACIONES}}} &
\multirow{2}{*}{\color{white}\scriptsize\textbf{\makecell{COMPETENCIAS}}} &
\multirow{2}{*}{\color{white}\scriptsize\textbf{\makecell{APRENDIZAJES}}} &
\multicolumn{3}{c|}{%
  \color{white}\makecell{\textbf{\footnotesize DESCRIPTORES DE DESEMPEÑO}\\[1pt]\textit{\scriptsize Metas del aprendizaje por periodo}}%
}\\
\cline{4-6}
\rowcolor{azulClaro}
\cellcolor{azulInst} & \cellcolor{azulInst} & \cellcolor{azulInst} &
\textbf{\scriptsize COGNITIVOS} &
\textbf{\scriptsize PROCEDIMENTALES} &
\textbf{\scriptsize ACTITUDINALES} \\
\hline
\endhead

% -- FILA 1: Naturaleza y evolución --
\textbf{Naturaleza y evolución de la tecnología} &
Reconocer las raíces ancestrales del pensamiento computacional en oficios tradicionales (tejido, cocina, agricultura). &
Identifica las 4 técnicas del PC y reconoce que las tejedoras ya programaban (S1–S3). &
Define algoritmo, secuencia, repetición, condicional y variable, y los reconoce en oficios cotidianos. &
Descompone un problema cotidiano (cepillarse, cocinar, tejer) en pasos algorítmicos verificables. &
Honra a las tejedoras del Valle como primeras programadoras y reconoce la dignidad de los saberes ancestrales. \\
\hline

% -- FILA 2: Apropiación y uso --
\textbf{Apropiación y uso de la tecnología} &
Programar en Scratch usando variables, condicionales y bucles con propósito claro. &
Construye programas en Scratch con personajes que responden al usuario (S8–S9). &
Conoce los bloques principales de Scratch: eventos, control, sensores, variables, operadores. &
Programa un juego o narrativa interactiva con al menos 2 personajes, 2 escenas, 2 variables, 2 condicionales y 2 bucles. &
Programa con paciencia y claridad: variables con nombre, bucles con propósito, condiciones sin duplicar. \\
\hline

% -- FILA 3: Solución de problemas --
\textbf{Solución de problemas con tecnología} &
Aplicar diagramas de flujo y depuración para resolver problemas algorítmicos. &
Elabora diagramas de flujo con figuras correctas (S4) y depura programas paso a paso (S8–S9). &
Reconoce las figuras del diagrama de flujo (óvalo/rectángulo/rombo/flecha) y su significado. &
Diseña un diagrama de flujo en cuaderno antes de programar; prueba el programa con usuarios externos. &
Acepta los errores como parte del proceso y depura con paciencia. \\
\hline

% -- FILA 4: Tecnología y sociedad --
\textbf{Tecnología y sociedad} &
Producir un programa con instructivo de uso para que otros lo entiendan y usen. &
Documenta su programa con instructivo de uso de 1 página (S10). &
Comprende que un programa debe ser entendible por otros, no solo por su autor. &
Redacta un instructivo de uso con 7 partes y lo prueba con un compañero externo al equipo. &
Reconoce que el código que escribe hoy se ejecuta automáticamente después; es responsable de lo que su programa hace. \\
\hline

% -- FILA 5: PC (verde) --
\rowcolor{verdePCclaro}
\makecell[l]{\textbf{\color{verdePC}Pensamiento}\\\textbf{\color{verdePC}Computacional}\\
\footnotesize Wing (2006)\\\footnotesize MEN Or. 2022\\\footnotesize\textit{Descomposición · Patrones · Abstracción · Algoritmos · Depuración}}\
 &
Aplicar las 4 técnicas centrales del PC (descomposición, patrones, abstracción, algoritmos) para diseñar y depurar programas en Scratch. &
Diseña algoritmos con secuencia, repetición y condicional (S1–S4); programa variables y bucles (S5–S8); depura y documenta (S9–S10). &
Define descomposición, reconocimiento de patrones, abstracción y diseño de algoritmos; los aplica en Scratch. &
Crea un programa en Scratch (2+ personajes, 2 escenas, 2 variables, 2 condicionales, 2 bucles); traza diagrama de flujo antes de programar. &
Acepta los errores como parte de la depuración; honra las tejedoras del Valle como primeras programadoras. \\
\hline

% -- FILA 6: SEL (violeta) --
\rowcolor{moradoSELclaro}
\makecell[l]{\textbf{\color{moradoSEL}Habilidades}\\\textbf{\color{moradoSEL}Socioemocionales}\\
\footnotesize Ley 2383/2024\\\footnotesize CASEL 2020\\\footnotesize\textit{Autogestión · Toma de decisiones}}\
 &
Desarrollar perseverancia y autogestión emocional frente a los errores de programación, aprendiendo a depurar con calma y a tomar decisiones metódicas ante problemas complejos. &
Reconoce la frustración como señal de aprendizaje y aplica estrategias de calma para depurar sin rendirse. &
Distingue entre error de sintaxis, error de lógica y error de diseño; identifica la estrategia adecuada para cada tipo. &
Aplica la rutina depuración-paso-a-paso: aislar el error, formular hipótesis, probar, registrar resultado. &
Persiste con ecuanimidad frente a errores difíciles; reconoce que depurar bien es tan valioso como programar bien. \\
\hline

\end{longtable}
\medskip
\noindent\textbf{Nota:} Las filas 5 (\textcolor{verdePC}{verde}) y 6 (\textcolor{moradoSEL}{violeta}) son componentes propios de este plan, adicionales a los cuatro componentes de la Guía No.~30 del MEN.

\subsubsection{Periodo 3 · Inteligencia artificial con criterio}
\noindent\begin{tabular}{ll}
\textbf{Grado:} Séptimo & \textbf{Periodo:} Tercero \\
\textbf{Proyecto integrador:} Consejeros digitales: IA con criterio para mi salón & \textbf{Anclaje:} El consejero del barrio (don Lucho, doña Mercedes, doña Esperanza) \\
\end{tabular}
\medskip

\begin{longtable}{|L{3.8cm}|L{3.3cm}|L{3.3cm}|L{2.5cm}|L{2.5cm}|L{2.5cm}|}
\caption{Estructura de malla curricular — G7° Periodo 3 · Inteligencia artificial con criterio}
\label{tab:g7p3}\\
\hline
\rowcolor{azulInst}
\multirow{2}{*}{\color{white}\scriptsize\textbf{\makecell{ESTÁNDARES, LINEAMIENTOS\\U ORIENTACIONES}}} &
\multirow{2}{*}{\color{white}\scriptsize\textbf{\makecell{COMPETENCIAS}}} &
\multirow{2}{*}{\color{white}\scriptsize\textbf{\makecell{APRENDIZAJES}}} &
\multicolumn{3}{c|}{%
  \color{white}\makecell{\textbf{\footnotesize DESCRIPTORES DE DESEMPEÑO}\\[1pt]\textit{\scriptsize Metas del aprendizaje por periodo}}%
}\\
\cline{4-6}
\rowcolor{azulClaro}
\cellcolor{azulInst} & \cellcolor{azulInst} & \cellcolor{azulInst} &
\textbf{\scriptsize COGNITIVOS} &
\textbf{\scriptsize PROCEDIMENTALES} &
\textbf{\scriptsize ACTITUDINALES} \\
\hline
\endfirsthead
\hline
\rowcolor{azulInst}
\multirow{2}{*}{\color{white}\scriptsize\textbf{\makecell{ESTÁNDARES, LINEAMIENTOS\\U ORIENTACIONES}}} &
\multirow{2}{*}{\color{white}\scriptsize\textbf{\makecell{COMPETENCIAS}}} &
\multirow{2}{*}{\color{white}\scriptsize\textbf{\makecell{APRENDIZAJES}}} &
\multicolumn{3}{c|}{%
  \color{white}\makecell{\textbf{\footnotesize DESCRIPTORES DE DESEMPEÑO}\\[1pt]\textit{\scriptsize Metas del aprendizaje por periodo}}%
}\\
\cline{4-6}
\rowcolor{azulClaro}
\cellcolor{azulInst} & \cellcolor{azulInst} & \cellcolor{azulInst} &
\textbf{\scriptsize COGNITIVOS} &
\textbf{\scriptsize PROCEDIMENTALES} &
\textbf{\scriptsize ACTITUDINALES} \\
\hline
\endhead

% -- FILA 1: Naturaleza y evolución --
\textbf{Naturaleza y evolución de la tecnología} &
Comprender qué es la inteligencia artificial, sus tipos y cómo aprende. &
Identifica IA en su vida diaria (S2), distingue IA estrecha de general (S3) y comprende el entrenamiento (S4). &
Define IA, LLM, dataset y entrenamiento; distingue aprendizaje supervisado y no supervisado. &
Identifica al menos 5 ejemplos de IA en su vida diaria (TikTok, Siri, filtros, etc.). &
Reconoce que el consejero del barrio aconsejaba pero no decidía; la IA hoy ocupa ese rol con sus límites. \\
\hline

% -- FILA 2: Apropiación y uso --
\textbf{Apropiación y uso de la tecnología} &
Formular prompts efectivos a modelos de lenguaje con criterio profesional. &
Aplica los 4 elementos del prompt (rol, contexto, tarea, formato) y lo itera (S6–S7). &
Conoce los componentes de un prompt efectivo y las técnicas de iteración. &
Construye una guía de 6 prompts útiles para grado 7°, probados en una IA real. &
Pide bien hoy para no editar mañana; el tiempo invertido en formular se recupera al multiplicar. \\
\hline

% -- FILA 3: Solución de problemas --
\textbf{Solución de problemas con tecnología} &
Verificar la información que produce una IA y detectar alucinaciones. &
Reconoce alucinaciones y deshonestidad académica con IA (S8); verifica datos contra fuentes externas. &
Identifica las alucinaciones como respuestas estadísticamente probables pero no verdaderas. &
Verifica al menos 2 datos generados por IA contra fuentes oficiales antes de incluirlos en un trabajo. &
Mantiene escepticismo saludable: la IA habla rápido y seguro, pero requiere verificación humana. \\
\hline

% -- FILA 4: Tecnología y sociedad --
\textbf{Tecnología y sociedad} &
Declarar honestamente el uso de IA en cualquier trabajo académico y aplicar la política institucional. &
Aplica la declaración honesta de uso de IA (S9) y redacta el decálogo ético del salón (S10). &
Conoce la política institucional de IA por grados y los 5 principios transversales no negociables. &
Incluye sección "Cómo trabajamos con IA" en sus entregables especificando modelo, partes, edición y verificación. &
Habita la infoesfera como productor responsable: la IA amplifica su voz, no la silencia. \\
\hline

% -- FILA 5: PC (verde) --
\rowcolor{verdePCclaro}
\makecell[l]{\textbf{\color{verdePC}Pensamiento}\\\textbf{\color{verdePC}Computacional}\\
\footnotesize Wing (2006)\\\footnotesize MEN Or. 2022\\\footnotesize\textit{Reconocimiento de patrones · Abstracción}}\
 &
Aplicar el pensamiento computacional para comprender cómo los modelos de IA reconocen patrones y abstraen información para producir respuestas. &
Identifica el proceso de entrenamiento de una IA como reconocimiento de patrones a gran escala (S2–S4); diseña prompts como instrucciones algorítmicas (S6–S7). &
Describe el aprendizaje automático como proceso de abstracción estadística de patrones; reconoce que un LLM es una función probabilística, no un razonador con conciencia. &
Diseña 6 prompts con estructura algorítmica (rol + contexto + tarea + formato + restricciones + verificación); itera y evalúa contra fuentes externas. &
Mantiene escepticismo saludable: verifica antes de aceptar; comprende que la IA no piensa sino que procesa patrones estadísticos. \\
\hline

% -- FILA 6: SEL (violeta) --
\rowcolor{moradoSELclaro}
\makecell[l]{\textbf{\color{moradoSEL}Habilidades}\\\textbf{\color{moradoSEL}Socioemocionales}\\
\footnotesize Ley 2383/2024\\\footnotesize CASEL 2020\\\footnotesize\textit{Autoconciencia · Toma de decisiones responsable}}\
 &
Reconocer cómo los algoritmos de IA pueden influir en las propias creencias y emociones, y desarrollar criterio autónomo para tomar decisiones informadas frente a contenido generado por IA. &
Identifica situaciones en las que una IA puede generar sesgos en sus opiniones o emociones (burbujas de filtro, contenido personalizado). &
Describe el mecanismo de personalización algorítmica y su impacto en la formación de creencias. &
Analiza su propia experiencia con plataformas algorítmicas e identifica 3 casos en que el algoritmo influyó en lo que creyó o sintió. &
Ejerce su autonomía crítica frente al contenido generado por IA; decide con criterio propio, no por lo que la IA produce más fácil. \\
\hline

\end{longtable}
\medskip
\noindent\textbf{Nota:} Las filas 5 (\textcolor{verdePC}{verde}) y 6 (\textcolor{moradoSEL}{violeta}) son componentes propios de este plan, adicionales a los cuatro componentes de la Guía No.~30 del MEN.

\subsection{Grado Octavo}

\subsubsection{Periodo 1 · Lógica computacional avanzada}
\noindent\begin{tabular}{ll}
\textbf{Grado:} Octavo & \textbf{Periodo:} Primero \\
\textbf{Proyecto integrador:} Decisiones lógicas: del fogón al algoritmo & \textbf{Anclaje:} El gesto de la cocina del Valle \\
\end{tabular}
\medskip

\begin{longtable}{|L{3.8cm}|L{3.3cm}|L{3.3cm}|L{2.5cm}|L{2.5cm}|L{2.5cm}|}
\caption{Estructura de malla curricular — G8° Periodo 1 · Lógica computacional avanzada}
\label{tab:g8p1}\\
\hline
\rowcolor{azulInst}
\multirow{2}{*}{\color{white}\scriptsize\textbf{\makecell{ESTÁNDARES, LINEAMIENTOS\\U ORIENTACIONES}}} &
\multirow{2}{*}{\color{white}\scriptsize\textbf{\makecell{COMPETENCIAS}}} &
\multirow{2}{*}{\color{white}\scriptsize\textbf{\makecell{APRENDIZAJES}}} &
\multicolumn{3}{c|}{%
  \color{white}\makecell{\textbf{\footnotesize DESCRIPTORES DE DESEMPEÑO}\\[1pt]\textit{\scriptsize Metas del aprendizaje por periodo}}%
}\\
\cline{4-6}
\rowcolor{azulClaro}
\cellcolor{azulInst} & \cellcolor{azulInst} & \cellcolor{azulInst} &
\textbf{\scriptsize COGNITIVOS} &
\textbf{\scriptsize PROCEDIMENTALES} &
\textbf{\scriptsize ACTITUDINALES} \\
\hline
\endfirsthead
\hline
\rowcolor{azulInst}
\multirow{2}{*}{\color{white}\scriptsize\textbf{\makecell{ESTÁNDARES, LINEAMIENTOS\\U ORIENTACIONES}}} &
\multirow{2}{*}{\color{white}\scriptsize\textbf{\makecell{COMPETENCIAS}}} &
\multirow{2}{*}{\color{white}\scriptsize\textbf{\makecell{APRENDIZAJES}}} &
\multicolumn{3}{c|}{%
  \color{white}\makecell{\textbf{\footnotesize DESCRIPTORES DE DESEMPEÑO}\\[1pt]\textit{\scriptsize Metas del aprendizaje por periodo}}%
}\\
\cline{4-6}
\rowcolor{azulClaro}
\cellcolor{azulInst} & \cellcolor{azulInst} & \cellcolor{azulInst} &
\textbf{\scriptsize COGNITIVOS} &
\textbf{\scriptsize PROCEDIMENTALES} &
\textbf{\scriptsize ACTITUDINALES} \\
\hline
\endhead

% -- FILA 1: Naturaleza y evolución --
\textbf{Naturaleza y evolución de la tecnología} &
Comprender la lógica booleana como base del razonamiento computacional moderno. &
Identifica operadores lógicos AND, OR, NOT y su origen en el álgebra de Boole. &
Define los 3 operadores lógicos principales y construye tablas de verdad. &
Resuelve ejercicios de lógica booleana con tablas de verdad verificadas. &
Honra el gesto ancestral de cuidado y verificación cruzada de la abuela cocinera. \\
\hline

% -- FILA 2: Apropiación y uso --
\textbf{Apropiación y uso de la tecnología} &
Programar condicionales compuestas y estructuras anidadas en Scratch o entornos similares. &
Construye programas con condicionales múltiples y anidadas para resolver problemas reales. &
Conoce las estructuras: si-entonces-sino, anidamiento, casos múltiples. &
Programa una simulación de decisiones cotidianas (vestirse según clima, planear menú según ingredientes). &
Acepta la complejidad lógica de la vida cotidiana sin simplificarla. \\
\hline

% -- FILA 3: Solución de problemas --
\textbf{Solución de problemas con tecnología} &
Aplicar lógica condicional compuesta para resolver problemas de la vida real. &
Diseña algoritmos con condiciones AND/OR para evitar errores típicos. &
Identifica casos donde AND vs OR cambia el resultado del algoritmo. &
Crea un algoritmo verificador de decisiones cotidianas con condicionales compuestas. &
Verifica cruzando condiciones antes de actuar, como la abuela frente al fogón. \\
\hline

% -- FILA 4: Tecnología y sociedad --
\textbf{Tecnología y sociedad} &
Aplicar lógica compuesta para evitar errores socialmente costosos en sistemas automatizados. &
Reflexiona sobre cómo algoritmos mal diseñados pueden excluir o perjudicar a grupos. &
Reconoce ejemplos de sistemas automatizados con consecuencias sociales (filtros bancarios, sistemas de admisión). &
Analiza un caso real de decisión algorítmica que tuvo impacto social y propone mejoras. &
Asume responsabilidad por los efectos de los algoritmos que diseña en otros. \\
\hline

% -- FILA 5: PC (verde) --
\rowcolor{verdePCclaro}
\makecell[l]{\textbf{\color{verdePC}Pensamiento}\\\textbf{\color{verdePC}Computacional}\\
\footnotesize Wing (2006)\\\footnotesize MEN Or. 2022\\\footnotesize\textit{Diseño de algoritmos · Abstracción · Depuración}}\
 &
Aplicar operadores lógicos booleanos (AND, OR, NOT) como fundamento del pensamiento computacional y razonamiento algorítmico avanzado. &
Construye tablas de verdad (S1–S3); diseña algoritmos con condicionales múltiples y anidados (S4–S7); depura programas lógicos (S8–S9). &
Define álgebra de Boole y sus 3 operadores; reconoce que todo procesamiento digital se reduce a operaciones lógicas binarias; abstrae problemas cotidianos como expresiones booleanas. &
Diseña algoritmos con condiciones AND/OR; crea una simulación de decisiones con lógica booleana verificable y traza la tabla de verdad. &
Verifica cruzando condiciones antes de decidir; asume responsabilidad por los efectos de los algoritmos que diseña. \\
\hline

% -- FILA 6: SEL (violeta) --
\rowcolor{moradoSELclaro}
\makecell[l]{\textbf{\color{moradoSEL}Habilidades}\\\textbf{\color{moradoSEL}Socioemocionales}\\
\footnotesize Ley 2383/2024\\\footnotesize CASEL 2020\\\footnotesize\textit{Autogestión · Conciencia social}}\
 &
Desarrollar regulación emocional frente a la frustración de la lógica compleja y conciencia social frente a los efectos de decisiones algorítmicas en comunidades reales. &
Reconoce la frustración que genera la lógica compleja como señal de aprendizaje y la regula con estrategias concretas; identifica el impacto humano de los errores algorítmicos. &
Distingue entre error lógico personal y consecuencia social de un algoritmo defectuoso; describe casos de impacto real. &
Aplica la técnica de descomposición del problema lógico cuando se bloquea; analiza un caso de algoritmo con impacto social negativo y propone mejoras. &
Regula su frustración con paciencia; asume responsabilidad social por los algoritmos que diseña, más allá de la corrección técnica. \\
\hline

\end{longtable}
\medskip
\noindent\textbf{Nota:} Las filas 5 (\textcolor{verdePC}{verde}) y 6 (\textcolor{moradoSEL}{violeta}) son componentes propios de este plan, adicionales a los cuatro componentes de la Guía No.~30 del MEN.

\subsubsection{Periodo 2 · Datos como tablas: pensar con estructura}
\noindent\begin{tabular}{ll}
\textbf{Grado:} Octavo & \textbf{Periodo:} Segundo \\
\textbf{Proyecto integrador:} Datos del barrio: ver lo que no se ve a simple vista & \textbf{Anclaje:} Decisiones del campesino en la finca \\
\end{tabular}
\medskip

\begin{longtable}{|L{3.8cm}|L{3.3cm}|L{3.3cm}|L{2.5cm}|L{2.5cm}|L{2.5cm}|}
\caption{Estructura de malla curricular — G8° Periodo 2 · Datos como tablas: pensar con estructura}
\label{tab:g8p2}\\
\hline
\rowcolor{azulInst}
\multirow{2}{*}{\color{white}\scriptsize\textbf{\makecell{ESTÁNDARES, LINEAMIENTOS\\U ORIENTACIONES}}} &
\multirow{2}{*}{\color{white}\scriptsize\textbf{\makecell{COMPETENCIAS}}} &
\multirow{2}{*}{\color{white}\scriptsize\textbf{\makecell{APRENDIZAJES}}} &
\multicolumn{3}{c|}{%
  \color{white}\makecell{\textbf{\footnotesize DESCRIPTORES DE DESEMPEÑO}\\[1pt]\textit{\scriptsize Metas del aprendizaje por periodo}}%
}\\
\cline{4-6}
\rowcolor{azulClaro}
\cellcolor{azulInst} & \cellcolor{azulInst} & \cellcolor{azulInst} &
\textbf{\scriptsize COGNITIVOS} &
\textbf{\scriptsize PROCEDIMENTALES} &
\textbf{\scriptsize ACTITUDINALES} \\
\hline
\endfirsthead
\hline
\rowcolor{azulInst}
\multirow{2}{*}{\color{white}\scriptsize\textbf{\makecell{ESTÁNDARES, LINEAMIENTOS\\U ORIENTACIONES}}} &
\multirow{2}{*}{\color{white}\scriptsize\textbf{\makecell{COMPETENCIAS}}} &
\multirow{2}{*}{\color{white}\scriptsize\textbf{\makecell{APRENDIZAJES}}} &
\multicolumn{3}{c|}{%
  \color{white}\makecell{\textbf{\footnotesize DESCRIPTORES DE DESEMPEÑO}\\[1pt]\textit{\scriptsize Metas del aprendizaje por periodo}}%
}\\
\cline{4-6}
\rowcolor{azulClaro}
\cellcolor{azulInst} & \cellcolor{azulInst} & \cellcolor{azulInst} &
\textbf{\scriptsize COGNITIVOS} &
\textbf{\scriptsize PROCEDIMENTALES} &
\textbf{\scriptsize ACTITUDINALES} \\
\hline
\endhead

% -- FILA 1: Naturaleza y evolución --
\textbf{Naturaleza y evolución de la tecnología} &
Comprender el dato como representación digital del mundo. &
Identifica tipos de datos: numérico, texto, fecha, booleano. &
Diferencia los tipos de datos y sus operaciones permitidas. &
Clasifica datos en categorías y los organiza en tablas. &
Reconoce que el campesino del Valle ya hacía análisis de datos sin Excel. \\
\hline

% -- FILA 2: Apropiación y uso --
\textbf{Apropiación y uso de la tecnología} &
Manejar hojas de cálculo (Excel/Google Sheets) con criterio profesional. &
Construye tablas con encabezados, fórmulas básicas (SUMA, PROMEDIO) y gráficos simples. &
Conoce las funciones básicas de una hoja de cálculo. &
Construye una tabla de 20+ filas con datos reales y aplica fórmulas y gráficos. &
Cuida la calidad del dato: completo, verificado, fechado. \\
\hline

% -- FILA 3: Solución de problemas --
\textbf{Solución de problemas con tecnología} &
Resolver problemas analizando datos estructurados con criterio estadístico básico. &
Calcula promedios, mínimos, máximos y elabora conclusiones a partir de datos. &
Identifica las estadísticas básicas y su interpretación. &
Analiza un conjunto de datos del barrio y obtiene conclusiones. &
Verifica los datos antes de concluir; evita generalizar desde muestras pequeñas. \\
\hline

% -- FILA 4: Tecnología y sociedad --
\textbf{Tecnología y sociedad} &
Reconocer el rol ético del dato en decisiones sociales y políticas. &
Reflexiona sobre cómo datos mal usados pueden discriminar o manipular. &
Reconoce sesgos comunes: muestra pequeña, no representativa, sesgo de confirmación. &
Analiza un caso de manipulación estadística y propone correcciones. &
Asume responsabilidad por los datos que produce y publica. \\
\hline

% -- FILA 5: PC (verde) --
\rowcolor{verdePCclaro}
\makecell[l]{\textbf{\color{verdePC}Pensamiento}\\\textbf{\color{verdePC}Computacional}\\
\footnotesize Wing (2006)\\\footnotesize MEN Or. 2022\\\footnotesize\textit{Abstracción · Reconocimiento de patrones · Descomposición}}\
 &
Aplicar la abstracción del PC para modelar información real como datos estructurados y reconocer patrones estadísticos en conjuntos de datos del entorno. &
Abstrae datos del barrio como registros en tabla (S1–S3); aplica fórmulas estadísticas como algoritmos de procesamiento (S4–S7); visualiza patrones (S8–S9). &
Describe los tipos de datos como abstracciones del mundo real; reconoce que la estructura de la tabla determina las operaciones algorítmicas posibles. &
Diseña el modelo de datos de su investigación de barrio: tablas con nombre de columna, tipo de dato y validación; verifica calidad antes de analizar. &
Cuida la calidad del dato como ética de la información; rechaza generalizar desde muestras pequeñas. \\
\hline

% -- FILA 6: SEL (violeta) --
\rowcolor{moradoSELclaro}
\makecell[l]{\textbf{\color{moradoSEL}Habilidades}\\\textbf{\color{moradoSEL}Socioemocionales}\\
\footnotesize Ley 2383/2024\\\footnotesize CASEL 2020\\\footnotesize\textit{Conciencia social · Toma de decisiones responsable}}\
 &
Reconocer la responsabilidad ética de quien maneja datos de personas reales de la comunidad y tomar decisiones responsables que respeten la dignidad y privacidad de los involucrados. &
Identifica los principios de privacidad y dignidad que deben guiar el manejo de datos de personas reales de su entorno. &
Describe la diferencia entre dato anónimo y dato personal; reconoce las implicaciones éticas del manejo irresponsable de datos comunitarios. &
Diseña su instrumento de recolección de datos con consentimiento implícito y anonimización antes de publicar resultados. &
Trata los datos de su comunidad con la misma dignidad con que trataría a las personas; no publica datos que puedan perjudicar a alguien. \\
\hline

\end{longtable}
\medskip
\noindent\textbf{Nota:} Las filas 5 (\textcolor{verdePC}{verde}) y 6 (\textcolor{moradoSEL}{violeta}) son componentes propios de este plan, adicionales a los cuatro componentes de la Guía No.~30 del MEN.

\subsubsection{Periodo 3 · Diseño visual y comunicación digital}
\noindent\begin{tabular}{ll}
\textbf{Grado:} Octavo & \textbf{Periodo:} Tercero \\
\textbf{Proyecto integrador:} Diseño visual con raíces: mi pieza editorial del Pacífico & \textbf{Anclaje:} Diseño visual del Pacífico colombiano \\
\end{tabular}
\medskip

\begin{longtable}{|L{3.8cm}|L{3.3cm}|L{3.3cm}|L{2.5cm}|L{2.5cm}|L{2.5cm}|}
\caption{Estructura de malla curricular — G8° Periodo 3 · Diseño visual y comunicación digital}
\label{tab:g8p3}\\
\hline
\rowcolor{azulInst}
\multirow{2}{*}{\color{white}\scriptsize\textbf{\makecell{ESTÁNDARES, LINEAMIENTOS\\U ORIENTACIONES}}} &
\multirow{2}{*}{\color{white}\scriptsize\textbf{\makecell{COMPETENCIAS}}} &
\multirow{2}{*}{\color{white}\scriptsize\textbf{\makecell{APRENDIZAJES}}} &
\multicolumn{3}{c|}{%
  \color{white}\makecell{\textbf{\footnotesize DESCRIPTORES DE DESEMPEÑO}\\[1pt]\textit{\scriptsize Metas del aprendizaje por periodo}}%
}\\
\cline{4-6}
\rowcolor{azulClaro}
\cellcolor{azulInst} & \cellcolor{azulInst} & \cellcolor{azulInst} &
\textbf{\scriptsize COGNITIVOS} &
\textbf{\scriptsize PROCEDIMENTALES} &
\textbf{\scriptsize ACTITUDINALES} \\
\hline
\endfirsthead
\hline
\rowcolor{azulInst}
\multirow{2}{*}{\color{white}\scriptsize\textbf{\makecell{ESTÁNDARES, LINEAMIENTOS\\U ORIENTACIONES}}} &
\multirow{2}{*}{\color{white}\scriptsize\textbf{\makecell{COMPETENCIAS}}} &
\multirow{2}{*}{\color{white}\scriptsize\textbf{\makecell{APRENDIZAJES}}} &
\multicolumn{3}{c|}{%
  \color{white}\makecell{\textbf{\footnotesize DESCRIPTORES DE DESEMPEÑO}\\[1pt]\textit{\scriptsize Metas del aprendizaje por periodo}}%
}\\
\cline{4-6}
\rowcolor{azulClaro}
\cellcolor{azulInst} & \cellcolor{azulInst} & \cellcolor{azulInst} &
\textbf{\scriptsize COGNITIVOS} &
\textbf{\scriptsize PROCEDIMENTALES} &
\textbf{\scriptsize ACTITUDINALES} \\
\hline
\endhead

% -- FILA 1: Naturaleza y evolución --
\textbf{Naturaleza y evolución de la tecnología} &
Reconocer la historia del diseño visual desde los petroglifos hasta la era digital. &
Identifica hitos del diseño gráfico y la influencia ancestral en la estética contemporánea. &
Conoce los principios de diseño: jerarquía, equilibrio, contraste, ritmo, unidad. &
Analiza piezas de diseño popular (almanaques, esquelas) y contemporáneas. &
Valora el diseño visual del Pacífico como tradición digna y vigente. \\
\hline

% -- FILA 2: Apropiación y uso --
\textbf{Apropiación y uso de la tecnología} &
Producir piezas visuales con herramientas digitales (Canva, Figma, Google Slides). &
Diseña piezas comunicativas (cartel, infografía, presentación) con criterio profesional. &
Identifica los principios de tipografía, color y layout. &
Produce una pieza visual de 1 página integrando texto, imagen e identidad visual. &
Cuida la accesibilidad: contrastes legibles, fuentes claras, jerarquía evidente. \\
\hline

% -- FILA 3: Solución de problemas --
\textbf{Solución de problemas con tecnología} &
Resolver problemas de comunicación visual con criterio ético y estético. &
Identifica cuándo un diseño confunde y cómo corregirlo. &
Reconoce errores comunes: exceso de fuentes, color saturado, jerarquía confusa. &
Revisa diseños propios y ajenos aplicando una lista de chequeo de 5 criterios. &
Acepta retroalimentación con humildad y mejora iterativamente. \\
\hline

% -- FILA 4: Tecnología y sociedad --
\textbf{Tecnología y sociedad} &
Producir piezas visuales que dialoguen con saberes y estéticas locales. &
Integra elementos del diseño visual del Pacífico en producciones propias. &
Reconoce los aportes estéticos de comunidades del Chocó, Cauca y Valle. &
Diseña una pieza visual final que integre al menos un elemento ancestral identificable. &
Honra a las comunidades del Pacífico citando sus aportes en cada pieza. \\
\hline

% -- FILA 5: PC (verde) --
\rowcolor{verdePCclaro}
\makecell[l]{\textbf{\color{verdePC}Pensamiento}\\\textbf{\color{verdePC}Computacional}\\
\footnotesize Wing (2006)\\\footnotesize MEN Or. 2022\\\footnotesize\textit{Descomposición · Reconocimiento de patrones · Evaluación}}\
 &
Aplicar la descomposición y reconocimiento de patrones del PC al análisis y producción de piezas visuales, identificando principios de diseño como patrones reutilizables. &
Descompone piezas de diseño en sus principios (jerarquía, contraste, ritmo, equilibrio, unidad) como unidades analizables (S1–S3); aplica patrones visuales identificados (S4–S8). &
Reconoce los principios de diseño como patrones reutilizables; describe la jerarquía visual como algoritmo de lectura. &
Analiza 3 piezas de diseño descomponiéndolas en sus patrones; produce su pieza visual con lista de chequeo sistemática de 5 criterios. &
Acepta la retroalimentación con humildad; mejora iterativamente; reconoce que el buen diseño sigue lógica observable. \\
\hline

% -- FILA 6: SEL (violeta) --
\rowcolor{moradoSELclaro}
\makecell[l]{\textbf{\color{moradoSEL}Habilidades}\\\textbf{\color{moradoSEL}Socioemocionales}\\
\footnotesize Ley 2383/2024\\\footnotesize CASEL 2020\\\footnotesize\textit{Autoconciencia · Habilidades relacionales}}\
 &
Desarrollar autoconciencia estética sobre el propio gusto y las raíces culturales, y habilidades relacionales para recibir y dar retroalimentación constructiva sobre el trabajo creativo. &
Distingue entre preferencia estética personal y criterio de diseño fundado; identifica cómo sus raíces culturales informan su gusto visual. &
Describe la diferencia entre opinión estética y criterio técnico; reconoce cómo la cultura del Pacífico aporta elementos visuales propios. &
Da retroalimentación de diseño usando los 5 criterios técnicos (no "no me gusta"); recibe retroalimentación sin defensividad y aplica mejoras. &
Honra las raíces culturales propias en sus decisiones estéticas; recibe y da retroalimentación con generosidad y criterio. \\
\hline

\end{longtable}
\medskip
\noindent\textbf{Nota:} Las filas 5 (\textcolor{verdePC}{verde}) y 6 (\textcolor{moradoSEL}{violeta}) son componentes propios de este plan, adicionales a los cuatro componentes de la Guía No.~30 del MEN.

\subsection{Grado Noveno}

\subsubsection{Periodo 1 · Técnica, ciencia y tecnología}
\noindent\begin{tabular}{ll}
\textbf{Grado:} Noveno & \textbf{Periodo:} Primero \\
\textbf{Proyecto integrador:} Manifiesto del técnico crítico & \textbf{Anclaje:} La técnica en la mano (azuela, hueso tallado) \\
\end{tabular}
\medskip

\begin{longtable}{|L{3.8cm}|L{3.3cm}|L{3.3cm}|L{2.5cm}|L{2.5cm}|L{2.5cm}|}
\caption{Estructura de malla curricular — G9° Periodo 1 · Técnica, ciencia y tecnología}
\label{tab:g9p1}\\
\hline
\rowcolor{azulInst}
\multirow{2}{*}{\color{white}\scriptsize\textbf{\makecell{ESTÁNDARES, LINEAMIENTOS\\U ORIENTACIONES}}} &
\multirow{2}{*}{\color{white}\scriptsize\textbf{\makecell{COMPETENCIAS}}} &
\multirow{2}{*}{\color{white}\scriptsize\textbf{\makecell{APRENDIZAJES}}} &
\multicolumn{3}{c|}{%
  \color{white}\makecell{\textbf{\footnotesize DESCRIPTORES DE DESEMPEÑO}\\[1pt]\textit{\scriptsize Metas del aprendizaje por periodo}}%
}\\
\cline{4-6}
\rowcolor{azulClaro}
\cellcolor{azulInst} & \cellcolor{azulInst} & \cellcolor{azulInst} &
\textbf{\scriptsize COGNITIVOS} &
\textbf{\scriptsize PROCEDIMENTALES} &
\textbf{\scriptsize ACTITUDINALES} \\
\hline
\endfirsthead
\hline
\rowcolor{azulInst}
\multirow{2}{*}{\color{white}\scriptsize\textbf{\makecell{ESTÁNDARES, LINEAMIENTOS\\U ORIENTACIONES}}} &
\multirow{2}{*}{\color{white}\scriptsize\textbf{\makecell{COMPETENCIAS}}} &
\multirow{2}{*}{\color{white}\scriptsize\textbf{\makecell{APRENDIZAJES}}} &
\multicolumn{3}{c|}{%
  \color{white}\makecell{\textbf{\footnotesize DESCRIPTORES DE DESEMPEÑO}\\[1pt]\textit{\scriptsize Metas del aprendizaje por periodo}}%
}\\
\cline{4-6}
\rowcolor{azulClaro}
\cellcolor{azulInst} & \cellcolor{azulInst} & \cellcolor{azulInst} &
\textbf{\scriptsize COGNITIVOS} &
\textbf{\scriptsize PROCEDIMENTALES} &
\textbf{\scriptsize ACTITUDINALES} \\
\hline
\endhead

% -- FILA 1: Naturaleza y evolución --
\textbf{Naturaleza y evolución de la tecnología} &
Comprender la evolución histórica de la técnica desde los primeros artefactos hasta el chip. &
Construye una línea de tiempo de la técnica desde el fuego hasta la era digital. &
Distingue técnica, ciencia y tecnología; reconoce hitos: fuego, rueda, imprenta, electricidad, internet. &
Investiga 5 hitos tecnológicos y los presenta en línea de tiempo visual. &
Reconoce que la técnica nació en la mano humana antes que en la fábrica. \\
\hline

% -- FILA 2: Apropiación y uso --
\textbf{Apropiación y uso de la tecnología} &
Usar herramientas tecnológicas con criterio histórico y crítico. &
Maneja herramientas digitales valorando lo que cada una tiene de oficios anteriores. &
Conoce el origen y propósito de cada herramienta que usa cotidianamente. &
Documenta el uso de 5 herramientas digitales con su contexto histórico. &
Cuida las herramientas como extensión del cuerpo, no como descartables. \\
\hline

% -- FILA 3: Solución de problemas --
\textbf{Solución de problemas con tecnología} &
Resolver problemas tecnológicos con perspectiva histórica y crítica. &
Compara soluciones técnicas antiguas y modernas al mismo problema. &
Identifica problemas que la humanidad ha resuelto varias veces de formas distintas. &
Analiza un problema actual y propone solución integrando tradición y tecnología. &
No idolatra lo nuevo ni desprecia lo antiguo; juzga por criterio. \\
\hline

% -- FILA 4: Tecnología y sociedad --
\textbf{Tecnología y sociedad} &
Construir manifiesto del técnico crítico responsable. &
Redacta un manifiesto personal sobre su rol como técnico/usuario responsable. &
Reconoce las dimensiones éticas, ambientales y sociales del trabajo técnico. &
Elabora un manifiesto de 1 página con sus principios técnicos personales. &
Asume posición crítica frente al consumismo tecnológico. \\
\hline

% -- FILA 5: PC (verde) --
\rowcolor{verdePCclaro}
\makecell[l]{\textbf{\color{verdePC}Pensamiento}\\\textbf{\color{verdePC}Computacional}\\
\footnotesize Wing (2006)\\\footnotesize MEN Or. 2022\\\footnotesize\textit{Reconocimiento de patrones · Modelado · Evaluación}}\
 &
Aplicar el modelado del PC para comparar soluciones técnicas de distintas épocas como algoritmos alternativos de resolución de problemas humanos. &
Construye línea de tiempo tecnológica como modelo comparativo (S1–S3); compara soluciones técnicas como alternativas algorítmicas al mismo problema humano. &
Reconoce que la técnica es un proceso sistemático de pasos para transformar materia; comprende la evolución tecnológica como refinamiento iterativo de algoritmos de solución. &
Analiza un problema tecnológico actual y propone 2 soluciones (tradicional y digital) como algoritmos comparables con criterios explícitos de eficiencia, costo y contexto. &
No idolatra lo nuevo ni desprecia lo antiguo; evalúa con criterio algorítmico comparativo; reconoce la dignidad de la técnica ancestral. \\
\hline

% -- FILA 6: SEL (violeta) --
\rowcolor{moradoSELclaro}
\makecell[l]{\textbf{\color{moradoSEL}Habilidades}\\\textbf{\color{moradoSEL}Socioemocionales}\\
\footnotesize Ley 2383/2024\\\footnotesize CASEL 2020\\\footnotesize\textit{Autoconciencia · Toma de decisiones responsable}}\
 &
Desarrollar autoconciencia sobre el propio rol como técnico o usuario de tecnología, y tomar decisiones responsables frente al consumismo tecnológico y el impacto ambiental de los artefactos. &
Identifica sus propias actitudes frente a lo nuevo (fascinación, rechazo, criterio) y las examina con honestidad. &
Describe el impacto ambiental del ciclo de vida de los dispositivos tecnológicos y las alternativas de consumo responsable. &
Redacta un manifiesto personal de 10 principios como técnico responsable, incluyendo al menos 2 principios sobre impacto ambiental. &
Asume posición crítica frente al consumismo tecnológico; aplica sus principios en decisiones concretas de uso y descarte. \\
\hline

\end{longtable}
\medskip
\noindent\textbf{Nota:} Las filas 5 (\textcolor{verdePC}{verde}) y 6 (\textcolor{moradoSEL}{violeta}) son componentes propios de este plan, adicionales a los cuatro componentes de la Guía No.~30 del MEN.

\subsubsection{Periodo 2 · Diseño editorial y maquetación digital}
\noindent\begin{tabular}{ll}
\textbf{Grado:} Noveno & \textbf{Periodo:} Segundo \\
\textbf{Proyecto integrador:} Revista escolar con diseño editorial & \textbf{Anclaje:} Pasquines, almanaques y esquelas del Valle \\
\end{tabular}
\medskip

\begin{longtable}{|L{3.8cm}|L{3.3cm}|L{3.3cm}|L{2.5cm}|L{2.5cm}|L{2.5cm}|}
\caption{Estructura de malla curricular — G9° Periodo 2 · Diseño editorial y maquetación digital}
\label{tab:g9p2}\\
\hline
\rowcolor{azulInst}
\multirow{2}{*}{\color{white}\scriptsize\textbf{\makecell{ESTÁNDARES, LINEAMIENTOS\\U ORIENTACIONES}}} &
\multirow{2}{*}{\color{white}\scriptsize\textbf{\makecell{COMPETENCIAS}}} &
\multirow{2}{*}{\color{white}\scriptsize\textbf{\makecell{APRENDIZAJES}}} &
\multicolumn{3}{c|}{%
  \color{white}\makecell{\textbf{\footnotesize DESCRIPTORES DE DESEMPEÑO}\\[1pt]\textit{\scriptsize Metas del aprendizaje por periodo}}%
}\\
\cline{4-6}
\rowcolor{azulClaro}
\cellcolor{azulInst} & \cellcolor{azulInst} & \cellcolor{azulInst} &
\textbf{\scriptsize COGNITIVOS} &
\textbf{\scriptsize PROCEDIMENTALES} &
\textbf{\scriptsize ACTITUDINALES} \\
\hline
\endfirsthead
\hline
\rowcolor{azulInst}
\multirow{2}{*}{\color{white}\scriptsize\textbf{\makecell{ESTÁNDARES, LINEAMIENTOS\\U ORIENTACIONES}}} &
\multirow{2}{*}{\color{white}\scriptsize\textbf{\makecell{COMPETENCIAS}}} &
\multirow{2}{*}{\color{white}\scriptsize\textbf{\makecell{APRENDIZAJES}}} &
\multicolumn{3}{c|}{%
  \color{white}\makecell{\textbf{\footnotesize DESCRIPTORES DE DESEMPEÑO}\\[1pt]\textit{\scriptsize Metas del aprendizaje por periodo}}%
}\\
\cline{4-6}
\rowcolor{azulClaro}
\cellcolor{azulInst} & \cellcolor{azulInst} & \cellcolor{azulInst} &
\textbf{\scriptsize COGNITIVOS} &
\textbf{\scriptsize PROCEDIMENTALES} &
\textbf{\scriptsize ACTITUDINALES} \\
\hline
\endhead

% -- FILA 1: Naturaleza y evolución --
\textbf{Naturaleza y evolución de la tecnología} &
Comprender la evolución del diseño editorial desde la imprenta de Gutenberg hasta el diseño web. &
Identifica hitos del diseño editorial y su raíz popular en pasquines y almanaques. &
Conoce los componentes editoriales: portada, contraportada, sumario, secciones, cierre. &
Investiga 3 revistas y describe su estructura editorial. &
Honra el diseño popular del Valle como tradición editorial vigente. \\
\hline

% -- FILA 2: Apropiación y uso --
\textbf{Apropiación y uso de la tecnología} &
Maquetar piezas editoriales con herramientas digitales (Canva, Figma, Adobe Express). &
Produce una revista de 8–12 páginas con maquetación coherente. &
Reconoce reglas de cuadrícula, tipografía editorial y jerarquía visual. &
Maqueta una revista escolar de 8+ páginas con secciones, imágenes y texto. &
Cuida cada detalle como el editor de tradición. \\
\hline

% -- FILA 3: Solución de problemas --
\textbf{Solución de problemas con tecnología} &
Resolver problemas de maquetación con criterio profesional. &
Detecta y corrige problemas editoriales (jerarquía confusa, fuentes mal combinadas, accesibilidad). &
Conoce los 10 errores comunes en diseño editorial. &
Revisa la revista propia y ajena aplicando lista de chequeo. &
Itera con paciencia; el primer borrador nunca es el final. \\
\hline

% -- FILA 4: Tecnología y sociedad --
\textbf{Tecnología y sociedad} &
Publicar piezas editoriales que aporten valor a la comunidad escolar. &
Distribuye su revista a la comunidad y registra retroalimentación. &
Reconoce el rol social de la prensa escolar y su capacidad de informar y formar. &
Sustenta la revista ante el salón y recoge feedback de mínimo 5 lectores. &
Asume responsabilidad por lo que publica como editor. \\
\hline

% -- FILA 5: PC (verde) --
\rowcolor{verdePCclaro}
\makecell[l]{\textbf{\color{verdePC}Pensamiento}\\\textbf{\color{verdePC}Computacional}\\
\footnotesize Wing (2006)\\\footnotesize MEN Or. 2022\\\footnotesize\textit{Descomposición · Abstracción · Depuración iterativa}}\
 &
Aplicar la descomposición y abstracción del PC para analizar la estructura editorial como algoritmo de comunicación con reglas, módulos y componentes interdependientes. &
Descompone la estructura de una revista (portada, sumario, secciones) como módulos interdependientes (S1–S3); aplica reglas de cuadrícula como algoritmo de layout (S4–S7). &
Describe la maquetación como proceso algorítmico con reglas tipográficas, de cuadrícula y de jerarquía. &
Diseña la arquitectura de su revista con esquema de módulos, reglas tipográficas y sistema de cuadrícula documentados; depura con checklist. &
Itera con paciencia; el primer borrador nunca es el final; aplica depuración editorial como el programador depura su código. \\
\hline

% -- FILA 6: SEL (violeta) --
\rowcolor{moradoSELclaro}
\makecell[l]{\textbf{\color{moradoSEL}Habilidades}\\\textbf{\color{moradoSEL}Socioemocionales}\\
\footnotesize Ley 2383/2024\\\footnotesize CASEL 2020\\\footnotesize\textit{Autogestión · Habilidades relacionales}}\
 &
Desarrollar perseverancia en el proceso iterativo de diseño editorial y habilidades relacionales para incorporar retroalimentación de lectores y colaborar en la producción de la revista escolar. &
Identifica sus patrones de respuesta ante el error de diseño (abandono, perfeccionismo, iteración) y desarrolla la iteración como hábito. &
Describe la retroalimentación del lector como dato de diseño, no como crítica personal. &
Incorpora mínimo 3 sugerencias de lectores reales en la versión final de la revista; documenta qué cambió y por qué. &
Recibe feedback de lectores con apertura y gratitud; itera la revista con disposición colaborativa. \\
\hline

\end{longtable}
\medskip
\noindent\textbf{Nota:} Las filas 5 (\textcolor{verdePC}{verde}) y 6 (\textcolor{moradoSEL}{violeta}) son componentes propios de este plan, adicionales a los cuatro componentes de la Guía No.~30 del MEN.

\subsubsection{Periodo 3 · Datos, estadística y visualización}
\noindent\begin{tabular}{ll}
\textbf{Grado:} Noveno & \textbf{Periodo:} Tercero \\
\textbf{Proyecto integrador:} Diagnóstico de mi colegio con datos & \textbf{Anclaje:} La libreta de la abuela (conteo de cosecha) \\
\end{tabular}
\medskip

\begin{longtable}{|L{3.8cm}|L{3.3cm}|L{3.3cm}|L{2.5cm}|L{2.5cm}|L{2.5cm}|}
\caption{Estructura de malla curricular — G9° Periodo 3 · Datos, estadística y visualización}
\label{tab:g9p3}\\
\hline
\rowcolor{azulInst}
\multirow{2}{*}{\color{white}\scriptsize\textbf{\makecell{ESTÁNDARES, LINEAMIENTOS\\U ORIENTACIONES}}} &
\multirow{2}{*}{\color{white}\scriptsize\textbf{\makecell{COMPETENCIAS}}} &
\multirow{2}{*}{\color{white}\scriptsize\textbf{\makecell{APRENDIZAJES}}} &
\multicolumn{3}{c|}{%
  \color{white}\makecell{\textbf{\footnotesize DESCRIPTORES DE DESEMPEÑO}\\[1pt]\textit{\scriptsize Metas del aprendizaje por periodo}}%
}\\
\cline{4-6}
\rowcolor{azulClaro}
\cellcolor{azulInst} & \cellcolor{azulInst} & \cellcolor{azulInst} &
\textbf{\scriptsize COGNITIVOS} &
\textbf{\scriptsize PROCEDIMENTALES} &
\textbf{\scriptsize ACTITUDINALES} \\
\hline
\endfirsthead
\hline
\rowcolor{azulInst}
\multirow{2}{*}{\color{white}\scriptsize\textbf{\makecell{ESTÁNDARES, LINEAMIENTOS\\U ORIENTACIONES}}} &
\multirow{2}{*}{\color{white}\scriptsize\textbf{\makecell{COMPETENCIAS}}} &
\multirow{2}{*}{\color{white}\scriptsize\textbf{\makecell{APRENDIZAJES}}} &
\multicolumn{3}{c|}{%
  \color{white}\makecell{\textbf{\footnotesize DESCRIPTORES DE DESEMPEÑO}\\[1pt]\textit{\scriptsize Metas del aprendizaje por periodo}}%
}\\
\cline{4-6}
\rowcolor{azulClaro}
\cellcolor{azulInst} & \cellcolor{azulInst} & \cellcolor{azulInst} &
\textbf{\scriptsize COGNITIVOS} &
\textbf{\scriptsize PROCEDIMENTALES} &
\textbf{\scriptsize ACTITUDINALES} \\
\hline
\endhead

% -- FILA 1: Naturaleza y evolución --
\textbf{Naturaleza y evolución de la tecnología} &
Comprender la estadística como ciencia de la incertidumbre y la visualización como lenguaje universal. &
Identifica los tipos de gráficos y cuándo usar cada uno. &
Define media, mediana, moda, rango, desviación; distingue gráficas de barras, líneas, sectores. &
Calcula medidas de tendencia central a partir de un conjunto de datos. &
Honra a la abuela cosechera como estadista ancestral. \\
\hline

% -- FILA 2: Apropiación y uso --
\textbf{Apropiación y uso de la tecnología} &
Producir visualizaciones honestas con hojas de cálculo y herramientas de BI básico. &
Diseña gráficos en Excel/Sheets con criterio comunicativo. &
Conoce los principios de la visualización honesta: ejes correctos, escalas justas, sin distorsión. &
Produce 3 visualizaciones distintas de un mismo dataset escolar. &
Diseña gráficos para informar, no para impresionar ni engañar. \\
\hline

% -- FILA 3: Solución de problemas --
\textbf{Solución de problemas con tecnología} &
Resolver problemas comunitarios con análisis de datos. &
Identifica un problema del colegio, recolecta datos y propone solución basada en evidencia. &
Reconoce los pasos del método de investigación con datos. &
Realiza encuesta a mínimo 30 estudiantes, analiza resultados y presenta hallazgos. &
Aplica rigor estadístico básico; reconoce los límites de su muestra. \\
\hline

% -- FILA 4: Tecnología y sociedad --
\textbf{Tecnología y sociedad} &
Comunicar resultados de análisis de datos a la comunidad escolar con honestidad. &
Presenta el diagnóstico del colegio a directivos y comunidad. &
Reconoce el rol del dato en la toma de decisiones democráticas. &
Elabora informe ciudadano de 6 páginas con datos, gráficos y propuestas. &
Honra la estadística doméstica de la abuela como ética de cuidado comunitario. \\
\hline

% -- FILA 5: PC (verde) --
\rowcolor{verdePCclaro}
\makecell[l]{\textbf{\color{verdePC}Pensamiento}\\\textbf{\color{verdePC}Computacional}\\
\footnotesize Wing (2006)\\\footnotesize MEN Or. 2022\\\footnotesize\textit{Algoritmos estadísticos · Abstracción · Evaluación}}\
 &
Aplicar algoritmos estadísticos y principios de abstracción del PC para analizar datos del entorno y comunicar hallazgos con honestidad y rigor. &
Aplica algoritmos estadísticos (promedio, mediana, moda, desviación) como funciones sobre conjuntos de datos (S1–S5); diseña visualizaciones según el patrón estadístico adecuado. &
Reconoce las medidas de tendencia central como algoritmos de síntesis de datos; identifica sesgos comunes como errores de abstracción (muestra no representativa, escala engañosa). &
Recolecta datos de mínimo 30 estudiantes, aplica los algoritmos estadísticos y diseña 3 visualizaciones distintas del mismo dataset; valida honestidad de sus gráficos. &
Aplica rigor estadístico; reconoce los límites de su muestra; asume responsabilidad por los datos que publica y la interpretación que propone. \\
\hline

% -- FILA 6: SEL (violeta) --
\rowcolor{moradoSELclaro}
\makecell[l]{\textbf{\color{moradoSEL}Habilidades}\\\textbf{\color{moradoSEL}Socioemocionales}\\
\footnotesize Ley 2383/2024\\\footnotesize CASEL 2020\\\footnotesize\textit{Conciencia social · Toma de decisiones responsable}}\
 &
Desarrollar conciencia social sobre el impacto que tienen los datos y las estadísticas en la comunidad escolar, y tomar decisiones responsables sobre cómo comunicar hallazgos sin manipular ni distorsionar. &
Identifica cómo una visualización deshonesta puede generar alarma, estigma o decisiones erróneas en la comunidad. &
Describe al menos 3 técnicas de distorsión estadística (escala truncada, muestra sesgada, correlación-causalidad) y sus consecuencias sociales. &
Revisa sus visualizaciones con un compañero usando la lista de chequeo de honestidad estadística antes de presentar a la comunidad. &
Asume la comunicación de datos como acto político y ético; presenta hallazgos con rigor, humildad y respeto por su comunidad. \\
\hline

\end{longtable}
\medskip
\noindent\textbf{Nota:} Las filas 5 (\textcolor{verdePC}{verde}) y 6 (\textcolor{moradoSEL}{violeta}) son componentes propios de este plan, adicionales a los cuatro componentes de la Guía No.~30 del MEN.

\subsection{Grado Décimo}

\subsubsection{Periodo 1 · Edición digital y autoría con IA responsable}
\noindent\begin{tabular}{ll}
\textbf{Grado:} Décimo & \textbf{Periodo:} Primero \\
\textbf{Proyecto integrador:} Mi libro digital: del oficio del editor al sello propio & \textbf{Anclaje:} El personaje silencioso del barrio (intermediario) \\
\end{tabular}
\medskip

\begin{longtable}{|L{3.8cm}|L{3.3cm}|L{3.3cm}|L{2.5cm}|L{2.5cm}|L{2.5cm}|}
\caption{Estructura de malla curricular — G10° Periodo 1 · Edición digital y autoría con IA responsable}
\label{tab:g10p1}\\
\hline
\rowcolor{azulInst}
\multirow{2}{*}{\color{white}\scriptsize\textbf{\makecell{ESTÁNDARES, LINEAMIENTOS\\U ORIENTACIONES}}} &
\multirow{2}{*}{\color{white}\scriptsize\textbf{\makecell{COMPETENCIAS}}} &
\multirow{2}{*}{\color{white}\scriptsize\textbf{\makecell{APRENDIZAJES}}} &
\multicolumn{3}{c|}{%
  \color{white}\makecell{\textbf{\footnotesize DESCRIPTORES DE DESEMPEÑO}\\[1pt]\textit{\scriptsize Metas del aprendizaje por periodo}}%
}\\
\cline{4-6}
\rowcolor{azulClaro}
\cellcolor{azulInst} & \cellcolor{azulInst} & \cellcolor{azulInst} &
\textbf{\scriptsize COGNITIVOS} &
\textbf{\scriptsize PROCEDIMENTALES} &
\textbf{\scriptsize ACTITUDINALES} \\
\hline
\endfirsthead
\hline
\rowcolor{azulInst}
\multirow{2}{*}{\color{white}\scriptsize\textbf{\makecell{ESTÁNDARES, LINEAMIENTOS\\U ORIENTACIONES}}} &
\multirow{2}{*}{\color{white}\scriptsize\textbf{\makecell{COMPETENCIAS}}} &
\multirow{2}{*}{\color{white}\scriptsize\textbf{\makecell{APRENDIZAJES}}} &
\multicolumn{3}{c|}{%
  \color{white}\makecell{\textbf{\footnotesize DESCRIPTORES DE DESEMPEÑO}\\[1pt]\textit{\scriptsize Metas del aprendizaje por periodo}}%
}\\
\cline{4-6}
\rowcolor{azulClaro}
\cellcolor{azulInst} & \cellcolor{azulInst} & \cellcolor{azulInst} &
\textbf{\scriptsize COGNITIVOS} &
\textbf{\scriptsize PROCEDIMENTALES} &
\textbf{\scriptsize ACTITUDINALES} \\
\hline
\endhead

% -- FILA 1: Naturaleza y evolución --
\textbf{Naturaleza y evolución de la tecnología} &
Comprender la evolución de la autoría editorial desde la imprenta hasta los libros con IA. &
Distingue el rol histórico del editor y los desafíos contemporáneos con IA generativa. &
Define qué hace y qué no hace un editor; reconoce el ciclo editorial completo. &
Analiza 3 libros editoriales y describe su estructura, voz y producción. &
Honra al personaje silencioso del barrio como intermediario cultural ancestral. \\
\hline

% -- FILA 2: Apropiación y uso --
\textbf{Apropiación y uso de la tecnología} &
Usar IA con criterio profesional para escritura editorial. &
Aplica las 5 partes del prompt profesional para texto y lo itera. &
Conoce las técnicas de prompting técnico para escritura: rol, contexto, restricciones, ejemplos, formato. &
Produce capítulos de libro con IA como asistente, editando con voz propia. &
La IA es asistente, no autor; la voz del libro es la del estudiante. \\
\hline

% -- FILA 3: Solución de problemas --
\textbf{Solución de problemas con tecnología} &
Resolver problemas de derechos de autor y autoría con IA según legislación colombiana. &
Aplica la Ley 23 de 1982 y las consideraciones éticas sobre obras derivadas con IA. &
Reconoce la diferencia entre obra original, obra derivada y obra generada por IA. &
Declara explícitamente la participación de IA en cada capítulo del libro. &
Sostiene la honestidad académica frente a la presión del atajo. \\
\hline

% -- FILA 4: Tecnología y sociedad --
\textbf{Tecnología y sociedad} &
Publicar un libro escolar con criterio profesional y compromiso ético. &
Compila, sustenta y difunde su libro de 80 páginas en la feria del libro escolar (S10). &
Reconoce el rol del autor responsable en la infoesfera contemporánea. &
Sustenta su libro públicamente y recoge retroalimentación crítica. &
Asume el rol de autor con responsabilidad y orgullo. \\
\hline

% -- FILA 5: PC (verde) --
\rowcolor{verdePCclaro}
\makecell[l]{\textbf{\color{verdePC}Pensamiento}\\\textbf{\color{verdePC}Computacional}\\
\footnotesize Wing (2006)\\\footnotesize MEN Or. 2022\\\footnotesize\textit{Diseño de algoritmos (prompting) · Depuración iterativa}}\
 &
Aplicar el PC para diseñar flujos de trabajo editoriales asistidos por IA, reconociendo el prompting como programación en lenguaje natural con estructura algorítmica. &
Diseña prompts como instrucciones algorítmicas (rol, contexto, tarea, restricciones) (S3–S5); itera el flujo editorial como proceso de depuración sistemática (S6–S8). &
Reconoce el prompt como instrucción algorítmica en lenguaje natural; comprende el flujo editorial asistido por IA como proceso con etapas, condiciones y revisiones documentadas. &
Diseña su flujo editorial para el libro digital (borrador IA / edición propia / verificación / versión final) como diagrama de proceso con decisiones y declaraciones documentadas. &
Asume la IA como asistente; la voz del libro es la suya; declara el uso con honestidad y responsabilidad ante sus lectores. \\
\hline

% -- FILA 6: SEL (violeta) --
\rowcolor{moradoSELclaro}
\makecell[l]{\textbf{\color{moradoSEL}Habilidades}\\\textbf{\color{moradoSEL}Socioemocionales}\\
\footnotesize Ley 2383/2024\\\footnotesize CASEL 2020\\\footnotesize\textit{Autoconciencia · Toma de decisiones responsable}}\
 &
Desarrollar autoconciencia sobre la propia voz autoral y tomar decisiones responsables frente al uso de IA, manteniendo integridad académica y reconociendo los límites éticos de la autoría asistida. &
Distingue con claridad cuáles partes del libro son voz propia y cuáles son asistencia de IA; evalúa si la proporción es éticamente adecuada. &
Describe los límites éticos de la autoría asistida por IA según la legislación colombiana y los principios de honestidad académica. &
Revisa cada capítulo de su libro con la pregunta "¿esta voz es mía?" y documenta las decisiones editoriales que tomó. &
Defiende su voz autoral como algo valioso e irremplazable; usa la IA para amplificarla, nunca para sustituirla. \\
\hline

\end{longtable}
\medskip
\noindent\textbf{Nota:} Las filas 5 (\textcolor{verdePC}{verde}) y 6 (\textcolor{moradoSEL}{violeta}) son componentes propios de este plan, adicionales a los cuatro componentes de la Guía No.~30 del MEN.

\subsubsection{Periodo 2 · Comunicación profesional con IA}
\noindent\begin{tabular}{ll}
\textbf{Grado:} Décimo & \textbf{Periodo:} Segundo \\
\textbf{Proyecto integrador:} Informe profesional del oficio & \textbf{Anclaje:} Los 4 oficios de la palabra (barrio Obrero) \\
\end{tabular}
\medskip

\begin{longtable}{|L{3.8cm}|L{3.3cm}|L{3.3cm}|L{2.5cm}|L{2.5cm}|L{2.5cm}|}
\caption{Estructura de malla curricular — G10° Periodo 2 · Comunicación profesional con IA}
\label{tab:g10p2}\\
\hline
\rowcolor{azulInst}
\multirow{2}{*}{\color{white}\scriptsize\textbf{\makecell{ESTÁNDARES, LINEAMIENTOS\\U ORIENTACIONES}}} &
\multirow{2}{*}{\color{white}\scriptsize\textbf{\makecell{COMPETENCIAS}}} &
\multirow{2}{*}{\color{white}\scriptsize\textbf{\makecell{APRENDIZAJES}}} &
\multicolumn{3}{c|}{%
  \color{white}\makecell{\textbf{\footnotesize DESCRIPTORES DE DESEMPEÑO}\\[1pt]\textit{\scriptsize Metas del aprendizaje por periodo}}%
}\\
\cline{4-6}
\rowcolor{azulClaro}
\cellcolor{azulInst} & \cellcolor{azulInst} & \cellcolor{azulInst} &
\textbf{\scriptsize COGNITIVOS} &
\textbf{\scriptsize PROCEDIMENTALES} &
\textbf{\scriptsize ACTITUDINALES} \\
\hline
\endfirsthead
\hline
\rowcolor{azulInst}
\multirow{2}{*}{\color{white}\scriptsize\textbf{\makecell{ESTÁNDARES, LINEAMIENTOS\\U ORIENTACIONES}}} &
\multirow{2}{*}{\color{white}\scriptsize\textbf{\makecell{COMPETENCIAS}}} &
\multirow{2}{*}{\color{white}\scriptsize\textbf{\makecell{APRENDIZAJES}}} &
\multicolumn{3}{c|}{%
  \color{white}\makecell{\textbf{\footnotesize DESCRIPTORES DE DESEMPEÑO}\\[1pt]\textit{\scriptsize Metas del aprendizaje por periodo}}%
}\\
\cline{4-6}
\rowcolor{azulClaro}
\cellcolor{azulInst} & \cellcolor{azulInst} & \cellcolor{azulInst} &
\textbf{\scriptsize COGNITIVOS} &
\textbf{\scriptsize PROCEDIMENTALES} &
\textbf{\scriptsize ACTITUDINALES} \\
\hline
\endhead

% -- FILA 1: Naturaleza y evolución --
\textbf{Naturaleza y evolución de la tecnología} &
Comprender la evolución de los géneros de texto profesional y sus convenciones. &
Identifica los 4 tipos de texto profesional y sus convenciones. &
Distingue informe técnico, informe comercial, ensayo y estudio de mercado. &
Analiza 4 textos profesionales y los clasifica por género. &
Honra a los 4 oficios de la palabra del barrio Obrero. \\
\hline

% -- FILA 2: Apropiación y uso --
\textbf{Apropiación y uso de la tecnología} &
Producir documentos profesionales con Google Docs/Word y IA responsable. &
Maneja estilos, índices, citas y referencias en documentos profesionales. &
Reconoce los componentes de un documento profesional formal. &
Produce un informe técnico de 10+ páginas con todos los componentes formales. &
Cuida cada detalle como el cuentista cuida la palabra empeñada. \\
\hline

% -- FILA 3: Solución de problemas --
\textbf{Solución de problemas con tecnología} &
Resolver problemas comunicativos en contextos laborales reales. &
Redacta correos profesionales, propuestas y reportes para audiencias reales. &
Identifica el tono, registro y estructura adecuados según audiencia. &
Redacta y envía 3 correos profesionales reales (a docentes, directivos o empresas). &
Sostiene profesionalismo aún en correos cortos. \\
\hline

% -- FILA 4: Tecnología y sociedad --
\textbf{Tecnología y sociedad} &
Aplicar la palabra profesional como herramienta de cambio social y empleabilidad. &
Sustenta el oficio de la palabra como ciudadanía activa (S10). &
Reconoce el valor profesional y social de saber escribir con criterio. &
Sustenta su informe ante audiencia real con criterio y argumentos. &
Honra la palabra empeñada como los 4 oficios del barrio Obrero. \\
\hline

% -- FILA 5: PC (verde) --
\rowcolor{verdePCclaro}
\makecell[l]{\textbf{\color{verdePC}Pensamiento}\\\textbf{\color{verdePC}Computacional}\\
\footnotesize Wing (2006)\\\footnotesize MEN Or. 2022\\\footnotesize\textit{Abstracción · Reconocimiento de patrones}}\
 &
Aplicar la abstracción y reconocimiento de patrones del PC para identificar las estructuras de género textual profesional como plantillas reutilizables y adaptables. &
Abstrae los componentes de 4 géneros profesionales como plantillas con patrones reutilizables (S1–S4); aplica y adapta las plantillas con voz propia (S5–S8). &
Reconoce los géneros textuales como abstracciones de patrones comunicativos con estructura, tono y propósito predefinidos; los describe como plantillas algorítmicas adaptables. &
Analiza 4 textos profesionales identificando sus patrones estructurales; produce un informe técnico siguiendo la plantilla identificada y adaptándola con voz propia verificable. &
Honra la palabra empeñada con la precisión del cuentista del barrio Obrero; sostiene profesionalismo aun en documentos cortos. \\
\hline

% -- FILA 6: SEL (violeta) --
\rowcolor{moradoSELclaro}
\makecell[l]{\textbf{\color{moradoSEL}Habilidades}\\\textbf{\color{moradoSEL}Socioemocionales}\\
\footnotesize Ley 2383/2024\\\footnotesize CASEL 2020\\\footnotesize\textit{Habilidades relacionales · Autogestión}}\
 &
Desarrollar habilidades relacionales para comunicarse con criterio profesional en contextos laborales y académicos, y autogestión emocional para mantener el tono adecuado aún bajo presión. &
Identifica situaciones de comunicación profesional que generan tensión emocional (correo de reclamo, informe con errores, comunicación con directivos) y practica respuestas asertivas. &
Describe la asertividad como equilibrio entre firmeza y cortesía; reconoce los efectos de la comunicación agresiva o pasiva en contextos profesionales. &
Redacta 3 correos profesionales en situaciones de tensión (reclamar, corregir, solicitar) usando registro asertivo y tono respetuoso. &
Mantiene el profesionalismo aún cuando el mensaje es difícil; honra la palabra empeñada también en la adversidad. \\
\hline

\end{longtable}
\medskip
\noindent\textbf{Nota:} Las filas 5 (\textcolor{verdePC}{verde}) y 6 (\textcolor{moradoSEL}{violeta}) son componentes propios de este plan, adicionales a los cuatro componentes de la Guía No.~30 del MEN.

\subsubsection{Periodo 3 · Microemprendimiento y gestión digital}
\noindent\begin{tabular}{ll}
\textbf{Grado:} Décimo & \textbf{Periodo:} Tercero \\
\textbf{Proyecto integrador:} Microemprendimiento escolar con plan completo & \textbf{Anclaje:} Tiendas y panaderías del centro (Cartago) \\
\end{tabular}
\medskip

\begin{longtable}{|L{3.8cm}|L{3.3cm}|L{3.3cm}|L{2.5cm}|L{2.5cm}|L{2.5cm}|}
\caption{Estructura de malla curricular — G10° Periodo 3 · Microemprendimiento y gestión digital}
\label{tab:g10p3}\\
\hline
\rowcolor{azulInst}
\multirow{2}{*}{\color{white}\scriptsize\textbf{\makecell{ESTÁNDARES, LINEAMIENTOS\\U ORIENTACIONES}}} &
\multirow{2}{*}{\color{white}\scriptsize\textbf{\makecell{COMPETENCIAS}}} &
\multirow{2}{*}{\color{white}\scriptsize\textbf{\makecell{APRENDIZAJES}}} &
\multicolumn{3}{c|}{%
  \color{white}\makecell{\textbf{\footnotesize DESCRIPTORES DE DESEMPEÑO}\\[1pt]\textit{\scriptsize Metas del aprendizaje por periodo}}%
}\\
\cline{4-6}
\rowcolor{azulClaro}
\cellcolor{azulInst} & \cellcolor{azulInst} & \cellcolor{azulInst} &
\textbf{\scriptsize COGNITIVOS} &
\textbf{\scriptsize PROCEDIMENTALES} &
\textbf{\scriptsize ACTITUDINALES} \\
\hline
\endfirsthead
\hline
\rowcolor{azulInst}
\multirow{2}{*}{\color{white}\scriptsize\textbf{\makecell{ESTÁNDARES, LINEAMIENTOS\\U ORIENTACIONES}}} &
\multirow{2}{*}{\color{white}\scriptsize\textbf{\makecell{COMPETENCIAS}}} &
\multirow{2}{*}{\color{white}\scriptsize\textbf{\makecell{APRENDIZAJES}}} &
\multicolumn{3}{c|}{%
  \color{white}\makecell{\textbf{\footnotesize DESCRIPTORES DE DESEMPEÑO}\\[1pt]\textit{\scriptsize Metas del aprendizaje por periodo}}%
}\\
\cline{4-6}
\rowcolor{azulClaro}
\cellcolor{azulInst} & \cellcolor{azulInst} & \cellcolor{azulInst} &
\textbf{\scriptsize COGNITIVOS} &
\textbf{\scriptsize PROCEDIMENTALES} &
\textbf{\scriptsize ACTITUDINALES} \\
\hline
\endhead

% -- FILA 1: Naturaleza y evolución --
\textbf{Naturaleza y evolución de la tecnología} &
Comprender la evolución de la gestión empresarial desde la libreta hasta el ERP. &
Identifica las herramientas digitales clave para la gestión moderna. &
Distingue ofimática, contabilidad básica, gestión de inventarios y CRM. &
Analiza la gestión digital de una microempresa real. &
Honra a las tiendas y panaderías como microempresas ancestrales. \\
\hline

% -- FILA 2: Apropiación y uso --
\textbf{Apropiación y uso de la tecnología} &
Manejar hojas de cálculo avanzadas para contabilidad personal y empresarial. &
Aplica fórmulas avanzadas, formato condicional y gráficos dinámicos. &
Conoce funciones BUSCARV, SI, SUMAR.SI, gráficos dinámicos. &
Construye libro de cuentas simulado con flujo de caja y punto de equilibrio. &
Cuida la calidad del dato contable: completo, fechado, verificado. \\
\hline

% -- FILA 3: Solución de problemas --
\textbf{Solución de problemas con tecnología} &
Resolver problemas de microemprendimiento con plan de negocio y Business Model Canvas. &
Construye un plan de negocio con IA para microempresa escolar. &
Identifica los 9 componentes del Business Model Canvas. &
Elabora plan de negocio de 15 páginas y BMC visual para su proyecto. &
Aplica el saber del comerciante del centro: planear y mantener la palabra. \\
\hline

% -- FILA 4: Tecnología y sociedad --
\textbf{Tecnología y sociedad} &
Publicar el microemprendimiento como ciudadanía económica responsable. &
Sustenta en feria de microempresas escolares (S10). &
Reconoce el rol social del emprendimiento responsable. &
Sustenta ante jurado real (docentes + invitados) y registra mejoras. &
Honra al comerciante tradicional como referente ético. \\
\hline

% -- FILA 5: PC (verde) --
\rowcolor{verdePCclaro}
\makecell[l]{\textbf{\color{verdePC}Pensamiento}\\\textbf{\color{verdePC}Computacional}\\
\footnotesize Wing (2006)\\\footnotesize MEN Or. 2022\\\footnotesize\textit{Modelado · Abstracción · Algoritmos de gestión}}\
 &
Aplicar el modelado del PC para diseñar el sistema de un microemprendimiento usando el Business Model Canvas como abstracción funcional de 9 componentes interrelacionados. &
Modela su microempresa usando el BMC como abstracción de 9 componentes con relaciones y flujos (S3–S6); usa hoja de cálculo como sistema de procesamiento algorítmico de datos financieros. &
Reconoce el Business Model Canvas como modelo de abstracción con 9 componentes, relaciones y flujos interdependientes. &
Diseña su BMC digital con 9 componentes completos; crea un libro de cuentas con fórmulas de flujo de caja, inventario y punto de equilibrio verificado. &
Aplica el saber del comerciante del centro: planear y mantener la palabra; cuida la calidad del dato contable como ética empresarial. \\
\hline

% -- FILA 6: SEL (violeta) --
\rowcolor{moradoSELclaro}
\makecell[l]{\textbf{\color{moradoSEL}Habilidades}\\\textbf{\color{moradoSEL}Socioemocionales}\\
\footnotesize Ley 2383/2024\\\footnotesize CASEL 2020\\\footnotesize\textit{Conciencia social · Toma de decisiones responsable}}\
 &
Desarrollar conciencia del impacto social del microemprendimiento en la comunidad y tomar decisiones empresariales responsables que honren la dignidad de clientes, empleados y del entorno. &
Identifica a los grupos impactados por su microemprendimiento (clientes, proveedores, entorno) y evalúa el impacto social de sus decisiones empresariales. &
Describe la responsabilidad social empresarial como componente del emprendimiento responsable; reconoce la diferencia entre ganancia cortoplacista y sostenibilidad a largo plazo. &
Incluye en su plan de negocio una sección de impacto social y ambiental con al menos 3 compromisos concretos y verificables. &
Emprende con la ética del comerciante del barrio: palabra empeñada, trato justo y compromiso con la comunidad como legado. \\
\hline

\end{longtable}
\medskip
\noindent\textbf{Nota:} Las filas 5 (\textcolor{verdePC}{verde}) y 6 (\textcolor{moradoSEL}{violeta}) son componentes propios de este plan, adicionales a los cuatro componentes de la Guía No.~30 del MEN.

\subsection{Grado Undécimo}

\subsubsection{Periodo 1 · Presencia digital empresarial con IA}
\noindent\begin{tabular}{ll}
\textbf{Grado:} Undécimo & \textbf{Periodo:} Primero \\
\textbf{Proyecto integrador:} Presencia digital empresarial con IA & \textbf{Anclaje:} Oficio y palabra (identidad y reputación digital) \\
\end{tabular}
\medskip

\begin{longtable}{|L{3.8cm}|L{3.3cm}|L{3.3cm}|L{2.5cm}|L{2.5cm}|L{2.5cm}|}
\caption{Estructura de malla curricular — G11° Periodo 1 · Presencia digital empresarial con IA}
\label{tab:g11p1}\\
\hline
\rowcolor{azulInst}
\multirow{2}{*}{\color{white}\scriptsize\textbf{\makecell{ESTÁNDARES, LINEAMIENTOS\\U ORIENTACIONES}}} &
\multirow{2}{*}{\color{white}\scriptsize\textbf{\makecell{COMPETENCIAS}}} &
\multirow{2}{*}{\color{white}\scriptsize\textbf{\makecell{APRENDIZAJES}}} &
\multicolumn{3}{c|}{%
  \color{white}\makecell{\textbf{\footnotesize DESCRIPTORES DE DESEMPEÑO}\\[1pt]\textit{\scriptsize Metas del aprendizaje por periodo}}%
}\\
\cline{4-6}
\rowcolor{azulClaro}
\cellcolor{azulInst} & \cellcolor{azulInst} & \cellcolor{azulInst} &
\textbf{\scriptsize COGNITIVOS} &
\textbf{\scriptsize PROCEDIMENTALES} &
\textbf{\scriptsize ACTITUDINALES} \\
\hline
\endfirsthead
\hline
\rowcolor{azulInst}
\multirow{2}{*}{\color{white}\scriptsize\textbf{\makecell{ESTÁNDARES, LINEAMIENTOS\\U ORIENTACIONES}}} &
\multirow{2}{*}{\color{white}\scriptsize\textbf{\makecell{COMPETENCIAS}}} &
\multirow{2}{*}{\color{white}\scriptsize\textbf{\makecell{APRENDIZAJES}}} &
\multicolumn{3}{c|}{%
  \color{white}\makecell{\textbf{\footnotesize DESCRIPTORES DE DESEMPEÑO}\\[1pt]\textit{\scriptsize Metas del aprendizaje por periodo}}%
}\\
\cline{4-6}
\rowcolor{azulClaro}
\cellcolor{azulInst} & \cellcolor{azulInst} & \cellcolor{azulInst} &
\textbf{\scriptsize COGNITIVOS} &
\textbf{\scriptsize PROCEDIMENTALES} &
\textbf{\scriptsize ACTITUDINALES} \\
\hline
\endhead

% -- FILA 1: Naturaleza y evolución --
\textbf{Naturaleza y evolución de la tecnología} &
Comprender la evolución de la presencia digital desde el currículum impreso al portafolio en línea. &
Identifica los componentes de una presencia digital profesional contemporánea. &
Distingue marca personal, identidad visual, sitio web personal, portafolio. &
Analiza 3 presencias digitales profesionales (uno excelente, uno regular, uno pobre). &
Honra el binomio ancestral oficio + palabra como anclaje de la reputación digital. \\
\hline

% -- FILA 2: Apropiación y uso --
\textbf{Apropiación y uso de la tecnología} &
Construir un sitio web personal profesional usando plantillas y IA. &
Maneja plantillas web (HTML/CSS responsive), copywriting con IA y SEO básico. &
Conoce los principios del diseño web responsivo, accesibilidad y SEO. &
Publica un sitio web personal en Netlify/Vercel con 4 secciones mínimas. &
Sostiene voz propia aún usando plantillas y IA: 30\% mínimo de intervención humana visible. \\
\hline

% -- FILA 3: Solución de problemas --
\textbf{Solución de problemas con tecnología} &
Diagnosticar y resolver problemas de presencia digital con criterio profesional. &
Diagnostica problemas (foto poco profesional, biografía sin voz, perfiles inconsistentes) y los corrige. &
Identifica los 10 errores comunes en presencia digital. &
Revisa su presencia digital con lista de chequeo y aplica correcciones. &
Acepta retroalimentación honesta sobre su marca personal. \\
\hline

% -- FILA 4: Tecnología y sociedad --
\textbf{Tecnología y sociedad} &
Lanzar su presencia digital al mundo con criterio ético y profesional. &
Compila perfiles, sitio web y métricas para su lanzamiento P1 (S10). &
Reconoce el rol de la presencia digital en empleabilidad y educación superior. &
Lanza presencia digital pública y registra métricas básicas (visitas, fuente). &
Asume reputación digital como compromiso de largo plazo. \\
\hline

% -- FILA 5: PC (verde) --
\rowcolor{verdePCclaro}
\makecell[l]{\textbf{\color{verdePC}Pensamiento}\\\textbf{\color{verdePC}Computacional}\\
\footnotesize Wing (2006)\\\footnotesize MEN Or. 2022\\\footnotesize\textit{Abstracción · Algoritmos (SEO) · Modelado}}\
 &
Aplicar el PC para diseñar una presencia digital sistemática, reconociendo los algoritmos de visibilidad (SEO) como sistemas de reglas con lógica propia y verificable. &
Analiza el algoritmo de SEO básico como sistema de reglas para visibilidad (S3–S5); diseña su presencia digital como arquitectura de información sistemática con jerarquía clara. &
Describe el SEO como algoritmo de clasificación basado en reglas de relevancia, autoridad y experiencia de usuario; reconoce la estructura web como jerarquía de información modelable. &
Diseña la arquitectura de su sitio web personal: mapa de secciones, jerarquía de contenidos, metadatos SEO y criterios de accesibilidad; documenta las decisiones con criterio algorítmico. &
Sostiene voz propia aún usando plantillas; asume la reputación digital como compromiso de largo plazo construido sistemáticamente. \\
\hline

% -- FILA 6: SEL (violeta) --
\rowcolor{moradoSELclaro}
\makecell[l]{\textbf{\color{moradoSEL}Habilidades}\\\textbf{\color{moradoSEL}Socioemocionales}\\
\footnotesize Ley 2383/2024\\\footnotesize CASEL 2020\\\footnotesize\textit{Autoconciencia · Toma de decisiones responsable}}\
 &
Desarrollar autoconciencia sobre la coherencia entre los valores personales y la presencia digital, y tomar decisiones responsables sobre la imagen pública que se proyecta al mundo profesional. &
Evalúa si su presencia digital actual refleja sus valores reales y la imagen profesional que desea proyectar a futuro. &
Describe las consecuencias de largo plazo de la huella digital; distingue entre imagen auténtica e imagen gestionada con integridad. &
Realiza una auditoría de su huella digital actual (redes, fotos, publicaciones) y diseña un plan de mejora con 5 acciones concretas. &
Construye su presencia digital con coherencia entre quien es y quien proyecta; asume la reputación digital como legado, no como estrategia. \\
\hline

\end{longtable}
\medskip
\noindent\textbf{Nota:} Las filas 5 (\textcolor{verdePC}{verde}) y 6 (\textcolor{moradoSEL}{violeta}) son componentes propios de este plan, adicionales a los cuatro componentes de la Guía No.~30 del MEN.

\subsubsection{Periodo 2 · Automatización y gestión de procesos}
\noindent\begin{tabular}{ll}
\textbf{Grado:} Undécimo & \textbf{Periodo:} Segundo \\
\textbf{Proyecto integrador:} Automatiza un proceso real & \textbf{Anclaje:} El maestro carpintero (lectura del proceso) \\
\end{tabular}
\medskip

\begin{longtable}{|L{3.8cm}|L{3.3cm}|L{3.3cm}|L{2.5cm}|L{2.5cm}|L{2.5cm}|}
\caption{Estructura de malla curricular — G11° Periodo 2 · Automatización y gestión de procesos}
\label{tab:g11p2}\\
\hline
\rowcolor{azulInst}
\multirow{2}{*}{\color{white}\scriptsize\textbf{\makecell{ESTÁNDARES, LINEAMIENTOS\\U ORIENTACIONES}}} &
\multirow{2}{*}{\color{white}\scriptsize\textbf{\makecell{COMPETENCIAS}}} &
\multirow{2}{*}{\color{white}\scriptsize\textbf{\makecell{APRENDIZAJES}}} &
\multicolumn{3}{c|}{%
  \color{white}\makecell{\textbf{\footnotesize DESCRIPTORES DE DESEMPEÑO}\\[1pt]\textit{\scriptsize Metas del aprendizaje por periodo}}%
}\\
\cline{4-6}
\rowcolor{azulClaro}
\cellcolor{azulInst} & \cellcolor{azulInst} & \cellcolor{azulInst} &
\textbf{\scriptsize COGNITIVOS} &
\textbf{\scriptsize PROCEDIMENTALES} &
\textbf{\scriptsize ACTITUDINALES} \\
\hline
\endfirsthead
\hline
\rowcolor{azulInst}
\multirow{2}{*}{\color{white}\scriptsize\textbf{\makecell{ESTÁNDARES, LINEAMIENTOS\\U ORIENTACIONES}}} &
\multirow{2}{*}{\color{white}\scriptsize\textbf{\makecell{COMPETENCIAS}}} &
\multirow{2}{*}{\color{white}\scriptsize\textbf{\makecell{APRENDIZAJES}}} &
\multicolumn{3}{c|}{%
  \color{white}\makecell{\textbf{\footnotesize DESCRIPTORES DE DESEMPEÑO}\\[1pt]\textit{\scriptsize Metas del aprendizaje por periodo}}%
}\\
\cline{4-6}
\rowcolor{azulClaro}
\cellcolor{azulInst} & \cellcolor{azulInst} & \cellcolor{azulInst} &
\textbf{\scriptsize COGNITIVOS} &
\textbf{\scriptsize PROCEDIMENTALES} &
\textbf{\scriptsize ACTITUDINALES} \\
\hline
\endhead

% -- FILA 1: Naturaleza y evolución --
\textbf{Naturaleza y evolución de la tecnología} &
Comprender la mejora continua (Kaizen) y los principios de la automatización. &
Identifica procesos repetitivos automatizables y sus indicadores (KPIs). &
Define BPMN, KPI, trigger, flujo y mejora continua. &
Mapea un proceso real con diagrama BPMN básico. &
Honra al maestro carpintero como lector del proceso. \\
\hline

% -- FILA 2: Apropiación y uso --
\textbf{Apropiación y uso de la tecnología} &
Construir automatizaciones de bajo código con triggers. &
Maneja Zapier, Make o n8n para automatizar tareas con disparadores. &
Conoce los componentes de una automatización: trigger + acción + filtros + transformación. &
Implementa una automatización real con mínimo 3 pasos. &
Cuida que la automatización no excluya casos legítimos. \\
\hline

% -- FILA 3: Solución de problemas --
\textbf{Solución de problemas con tecnología} &
Diagnosticar y resolver cuellos de botella en procesos del entorno. &
Identifica dónde se pierde tiempo en un proceso y propone solución. &
Reconoce las 7 mudas (desperdicios) del pensamiento Lean. &
Analiza un proceso real y diseña 3 mejoras priorizadas. &
Aplica Kaizen: mejora pequeña, continua, sostenida. \\
\hline

% -- FILA 4: Tecnología y sociedad --
\textbf{Tecnología y sociedad} &
Aplicar IA en mensajería responsable (chatbots con criterio). &
Construye un chatbot sencillo con flujo de conversación claro. &
Reconoce las consideraciones éticas de mensajería con IA: claridad sobre qué es IA, privacidad. &
Implementa un chatbot real para un caso concreto del colegio o la casa. &
Honra la palabra: si es chatbot, lo declara abiertamente al usuario. \\
\hline

% -- FILA 5: PC (verde) --
\rowcolor{verdePCclaro}
\makecell[l]{\textbf{\color{verdePC}Pensamiento}\\\textbf{\color{verdePC}Computacional}\\
\footnotesize Wing (2006)\\\footnotesize MEN Or. 2022\\\footnotesize\textit{Algoritmos · Descomposición de procesos · Abstracción (BPMN)}}\
 &
Aplicar el PC para modelar procesos con BPMN como diagramas de flujo y diseñar automatizaciones como programación de bajo código con triggers, condiciones y acciones. &
Modela un proceso real con BPMN básico (S1–S3); implementa automatización en Zapier/Make como algoritmo de trigger-acción-filtro (S4–S7); aplica Kaizen como depuración iterativa. &
Reconoce BPMN como lenguaje de modelado algorítmico de procesos; comprende la automatización como programa con trigger (evento), condición (filtro) y acción (resultado) verificable. &
Mapea un proceso del colegio/casa con diagrama BPMN; implementa una automatización de mínimo 3 pasos con trigger, condición y acción documentados y verificados con usuarios. &
Honra al maestro carpintero: lee el proceso antes de intervenir; mejora pequeña, continua y sostenida; la automatización no excluye casos legítimos. \\
\hline

% -- FILA 6: SEL (violeta) --
\rowcolor{moradoSELclaro}
\makecell[l]{\textbf{\color{moradoSEL}Habilidades}\\\textbf{\color{moradoSEL}Socioemocionales}\\
\footnotesize Ley 2383/2024\\\footnotesize CASEL 2020\\\footnotesize\textit{Conciencia social · Toma de decisiones responsable}}\
 &
Desarrollar conciencia social sobre el impacto de la automatización en el empleo y la vida de personas reales, y tomar decisiones responsables que preserven la dignidad y la inclusión frente a los sistemas automatizados. &
Identifica los riesgos de desplazamiento laboral y exclusión que puede generar una automatización mal diseñada, y los evalúa antes de implementar. &
Describe los efectos históricos de la automatización en el mercado laboral; distingue entre automatización que libera y automatización que excluye. &
Incluye en el diseño de su automatización un análisis de impacto humano: quién se beneficia, quién podría quedar excluido y cómo mitigarlo. &
Diseña automatizaciones con criterio humano: la eficiencia nunca justifica la exclusión; asume la responsabilidad de los sistemas que pone en marcha. \\
\hline

\end{longtable}
\medskip
\noindent\textbf{Nota:} Las filas 5 (\textcolor{verdePC}{verde}) y 6 (\textcolor{moradoSEL}{violeta}) son componentes propios de este plan, adicionales a los cuatro componentes de la Guía No.~30 del MEN.

\subsubsection{Periodo 3 · Microempresa con MVP, datos y ética}
\noindent\begin{tabular}{ll}
\textbf{Grado:} Undécimo & \textbf{Periodo:} Tercero \\
\textbf{Proyecto integrador:} Microempresa con MVP, datos y ética & \textbf{Anclaje:} Oficios que escuchan al pueblo (validación de mercado ancestral) \\
\end{tabular}
\medskip

\begin{longtable}{|L{3.8cm}|L{3.3cm}|L{3.3cm}|L{2.5cm}|L{2.5cm}|L{2.5cm}|}
\caption{Estructura de malla curricular — G11° Periodo 3 · Microempresa con MVP, datos y ética}
\label{tab:g11p3}\\
\hline
\rowcolor{azulInst}
\multirow{2}{*}{\color{white}\scriptsize\textbf{\makecell{ESTÁNDARES, LINEAMIENTOS\\U ORIENTACIONES}}} &
\multirow{2}{*}{\color{white}\scriptsize\textbf{\makecell{COMPETENCIAS}}} &
\multirow{2}{*}{\color{white}\scriptsize\textbf{\makecell{APRENDIZAJES}}} &
\multicolumn{3}{c|}{%
  \color{white}\makecell{\textbf{\footnotesize DESCRIPTORES DE DESEMPEÑO}\\[1pt]\textit{\scriptsize Metas del aprendizaje por periodo}}%
}\\
\cline{4-6}
\rowcolor{azulClaro}
\cellcolor{azulInst} & \cellcolor{azulInst} & \cellcolor{azulInst} &
\textbf{\scriptsize COGNITIVOS} &
\textbf{\scriptsize PROCEDIMENTALES} &
\textbf{\scriptsize ACTITUDINALES} \\
\hline
\endfirsthead
\hline
\rowcolor{azulInst}
\multirow{2}{*}{\color{white}\scriptsize\textbf{\makecell{ESTÁNDARES, LINEAMIENTOS\\U ORIENTACIONES}}} &
\multirow{2}{*}{\color{white}\scriptsize\textbf{\makecell{COMPETENCIAS}}} &
\multirow{2}{*}{\color{white}\scriptsize\textbf{\makecell{APRENDIZAJES}}} &
\multicolumn{3}{c|}{%
  \color{white}\makecell{\textbf{\footnotesize DESCRIPTORES DE DESEMPEÑO}\\[1pt]\textit{\scriptsize Metas del aprendizaje por periodo}}%
}\\
\cline{4-6}
\rowcolor{azulClaro}
\cellcolor{azulInst} & \cellcolor{azulInst} & \cellcolor{azulInst} &
\textbf{\scriptsize COGNITIVOS} &
\textbf{\scriptsize PROCEDIMENTALES} &
\textbf{\scriptsize ACTITUDINALES} \\
\hline
\endhead

% -- FILA 1: Naturaleza y evolución --
\textbf{Naturaleza y evolución de la tecnología} &
Comprender el ciclo emprendedor moderno y su raíz en escuchar al pueblo. &
Identifica los 10 pasos del ciclo emprendedor desde la idea al lanzamiento. &
Define MVP, validación, pivote, traction, pitch. &
Documenta su recorrido emprendedor con bitácora. &
Honra a los oficios del Valle que nacían de escuchar al pueblo. \\
\hline

% -- FILA 2: Apropiación y uso --
\textbf{Apropiación y uso de la tecnología} &
Construir un MVP digital y recolectar datos de validación reales. &
Implementa un MVP en plataforma de bajo código y recolecta datos de mínimo 20 usuarios. &
Conoce las plataformas no-code para MVPs (Glide, Bubble, sitios web). &
Implementa un MVP funcional y recolecta datos de mínimo 20 usuarios reales. &
Acepta el MVP como imperfecto pero útil para aprender. \\
\hline

% -- FILA 3: Solución de problemas --
\textbf{Solución de problemas con tecnología} &
Resolver problemas reales de microemprendimiento con datos, presupuesto y ética. &
Construye plan completo: ideación, validación, modelo, MVP, datos, pitch, presupuesto, ética, versión final. &
Reconoce el ciclo completo del emprendimiento responsable. &
Entrega plan completo de 20+ páginas con todas las secciones. &
Sostiene la ética en cada decisión (palabra empeñada vs atajo deshonesto). \\
\hline

% -- FILA 4: Tecnología y sociedad --
\textbf{Tecnología y sociedad} &
Sustentar el microemprendimiento ante una comunidad real con compromiso ético. &
Sustenta ante jurado real (docentes + emprendedores) con declaración ética de uso de IA (S10). &
Reconoce su rol como emprendedor responsable en la economía y la infoesfera. &
Sustenta públicamente y registra compromiso de mejora a 6 meses. &
Honra al pueblo como destinatario y validador del oficio emprendedor. \\
\hline

% -- FILA 5: PC (verde) --
\rowcolor{verdePCclaro}
\makecell[l]{\textbf{\color{verdePC}Pensamiento}\\\textbf{\color{verdePC}Computacional}\\
\footnotesize Wing (2006)\\\footnotesize MEN Or. 2022\\\footnotesize\textit{Prototipado iterativo (MVP) · Depuración con datos · Modelado}}\
 &
Aplicar el PC para diseñar un ciclo emprendedor iterativo con MVP como prototipo mínimo verificable, validación de datos y depuración basada en evidencia real. &
Diseña el MVP como prototipo mínimo verificable (S3–S5); valida con datos de mínimo 20 usuarios reales; itera y pivota basado en evidencia (S6–S8). &
Reconoce el ciclo emprendedor (idea / modelo / MVP / validación / pivote / escala) como algoritmo iterativo de depuración; comprende el MVP como prototipo mínimo con criterio de verificación definido. &
Implementa un MVP funcional en plataforma no-code; recolecta datos de mínimo 20 usuarios; analiza resultados y documenta 3 pivotes o mejoras con evidencia verificable. &
Acepta el MVP como imperfecto pero útil para aprender; sostiene la ética en cada decisión; honra al pueblo como destinatario y validador real del oficio emprendedor. \\
\hline

% -- FILA 6: SEL (violeta) --
\rowcolor{moradoSELclaro}
\makecell[l]{\textbf{\color{moradoSEL}Habilidades}\\\textbf{\color{moradoSEL}Socioemocionales}\\
\footnotesize Ley 2383/2024\\\footnotesize CASEL 2020\\\footnotesize\textit{Autogestión · Conciencia social}}\
 &
Desarrollar resiliencia emprendedora para afrontar fracasos y pivotes con ecuanimidad, y conciencia social para emprender con responsabilidad hacia la comunidad y el entorno como legado del oficio ancestral. &
Identifica sus respuestas emocionales frente al fracaso del MVP o el rechazo de usuarios, y desarrolla estrategias de resiliencia y reencuadre productivo. &
Describe la resiliencia emprendedora como capacidad de aprender del fracaso sin perder la orientación al propósito; reconoce el pivote como decisión estratégica, no como derrota. &
Documenta en su bitácora al menos 2 momentos de fracaso o rechazo, la emoción que generó y la decisión que tomó a partir de ellos. &
Emprende con ecuanimidad y propósito social; honra al pueblo como destinatario del oficio; deja el legado del emprendedor que escucha antes de vender. \\
\hline

\end{longtable}
\medskip
\noindent\textbf{Nota:} Las filas 5 (\textcolor{verdePC}{verde}) y 6 (\textcolor{moradoSEL}{violeta}) son componentes propios de este plan, adicionales a los cuatro componentes de la Guía No.~30 del MEN.

